\section{The memory of several sheets}
\label{sec:peps_sheet_stacks}
In the protocol of~\cite[Proposition~7.1]{OpenAI2026PolynomialPEPS} the memory
holds the main sheet, the active auxiliary sheets, the frozen auxiliary sheets
and the copies of $\Omega$ not yet in use. Every copy of $\Omega$, including the
main copy at the root, is initialized privately at its future owner, and the
reference vector at every stage is the tensor product of the reference vectors
of the encoded frames of the sheets, every factor of norm at most one
(\texttt{06-geometry.tex}, lines 653--681). This section collects the
bookkeeping of this memory: the initial private maps, the reorderings of the
sheets, the changes that leave a frame unchanged, and the norm of the reference
vector. The initialization assumes that the product of the reference vectors of the
frames of each party has norm one; for the raw frames of the initial memory, which have
no holes, this is the normalization of $\Omega$.

\begin{definition}[The memory of several sheets]
    \label{def:peps_sheet_stack}
    \lean{TNLean.PEPS.EncodedFrame.stackRegs,
        TNLean.PEPS.EncodedFrame.SheetStack,
        TNLean.PEPS.EncodedFrame.Frame.regRef,
        TNLean.PEPS.EncodedFrame.stackVec,
        TNLean.PEPS.EncodedFrame.framesVec}
    \leanok
    \uses{def:peps_frame_registers}
    For encoded frames $\mathcal F_1,\dots,\mathcal F_m$ (Definition~\ref{def:peps_frame_registers}) and a list $\ell$ of further registers with vector $\xi$, the sheet memory is the ordered tensor product of the registers of $\mathcal F_1,\dots,\mathcal F_m$ followed by $\ell$, with reference vector $\Omega_{\mathcal F_1}\otimes\cdots\otimes\Omega_{\mathcal F_m}\otimes\xi$.
\end{definition}

\begin{theorem}[Reordering the sheets]
    \label{thm:peps_sheet_reordering}
    \lean{TNLean.PEPS.EncodedFrame.stackRegs_append,
        TNLean.PEPS.EncodedFrame.stackVec_append,
        TNLean.PEPS.EncodedFrame.exists_reordering_stackVec,
        TNLean.PEPS.EncodedFrame.exists_originalCircuit_stackVec_of_perm}
    \leanok
    \uses{def:peps_sheet_stack}
    For any permutation $\mathcal F'_1,\dots,\mathcal F'_m$ of $\mathcal F_1,\dots,\mathcal F_m$, a word of tensor-factor exchanges (with $0$ nonprivate gates as an original circuit) maps the sheet memory of $\mathcal F_1,\dots,\mathcal F_m,\ell$ to that of $\mathcal F'_1,\dots,\mathcal F'_m,\ell$ and sends
    \begin{align}
        \Omega_{\mathcal F_1}\otimes\cdots\otimes\Omega_{\mathcal F_m}\otimes\xi
         & \mapsto \Omega_{\mathcal F'_1}\otimes\cdots\otimes\Omega_{\mathcal F'_m}\otimes\xi \quad (\forall \xi). \notag
    \end{align}
    In particular, an auxiliary sheet can be moved adjacent to the main sheet before a lens exchange.
\end{theorem}

\begin{proof}\leanok
    Induct on the permutation: prepending a common frame tensored on the left frames the word by its registers, and two adjacent frames are swapped by exchanging their register blocks.
\end{proof}

\begin{theorem}[The initial private maps]
    \label{thm:peps_sheet_initialization}
    \lean{TNLean.PEPS.EncodedFrame.framesSpace,
        TNLean.PEPS.EncodedFrame.exists_originalCircuit_initial}
    \leanok
    \uses{def:peps_sheet_stack,thm:lattice_stack_initialization}
    Suppose each frame in $\mathcal F_1,\dots,\mathcal F_m$ has all its raw and tag registers owned by a single party, and each party $p$ holds as its initial private memory the registers of its frames in list order with unit-norm product reference vector $\bigotimes_{\mathcal F_i\text{ at }p}\Omega_{\mathcal F_i}$. Then one private map per party followed by a sheet reordering (with $0$ nonprivate gates) maps $\bigotimes_p\bigotimes_{\mathcal F_i\text{ at }p}\Omega_{\mathcal F_i}$ exactly to $\Omega_{\mathcal F_1}\otimes\cdots\otimes\Omega_{\mathcal F_m}$ (\texttt{06-geometry.tex}, lines 658--662).
\end{theorem}

\begin{proof}\leanok
    \uses{thm:peps_sheet_reordering}
    Each private map identifies a party's memory with the tensor product of its frames' memories. Concatenating the per-party frame lists yields a permutation of $\mathcal F_1,\dots,\mathcal F_m$, which Theorem~\ref{thm:peps_sheet_reordering} reorders.
\end{proof}

\begin{theorem}[The initial step from raw copies]
    \label{thm:peps_sheet_raw_initialization}
    \lean{TNLean.PEPS.EncodedFrame.Frame.norm_refVec_of_holes_eq_nil,
        TNLean.PEPS.EncodedFrame.norm_framesVec_of_holes_eq_nil,
        TNLean.PEPS.EncodedFrame.exists_originalCircuit_initial_of_holes_eq_nil}
    \leanok
    \uses{thm:peps_sheet_initialization}
    Any hole-free frame $\mathcal F$ satisfies $\lVert\Omega_{\mathcal F}\rVert=\lVert\Omega\rVert$. If $\lVert\Omega\rVert=1$ and every frame in the list is hole-free and held entirely by one party (the main copy at the root and auxiliary copies at their future owners), then each party's product reference vector has unit norm and Theorem~\ref{thm:peps_sheet_initialization} applies (\texttt{06-geometry.tex}, lines 658--667).
\end{theorem}

\begin{proof}\leanok
    \uses{thm:peps_sheet_initialization, thm:peps_frame_norm}
    With no holes, the hole-projector product is $\id$, so $\lVert\Omega_{\mathcal F}\rVert=\lVert\Omega\rVert=1$ (Theorem~\ref{thm:peps_frame_norm}) and tensor products preserve unit norm.
\end{proof}

\begin{theorem}[Changes that leave a frame unchanged]
    \label{thm:peps_identity_change}
    \lean{TNLean.PEPS.EncodedFrame.exists_originalCircuit_of_frame_eq,
        TNLean.PEPS.EncodedFrame.Frame.changeOwner_empty,
        TNLean.PEPS.Approximation.emptyTarget_step}
    \leanok
    \uses{def:peps_sheet_stack}
    Changing the owner of $\varnothing$ leaves a frame unchanged: $\mathcal F[\varnothing\mapsto Q]=\mathcal F$. A transition between equal frames (such as an empty-target birth/death or an unchanged-radius hole resize) is implemented by the empty circuit, which has $0$ nonprivate gates and fixes $\Omega_{\mathcal F}$ (\texttt{06-geometry.tex}, lines 432--438 and 529--541).
\end{theorem}

\begin{proof}\leanok
    Immediate from Definition~\ref{def:peps_sheet_stack}.
\end{proof}

\begin{theorem}[The norm of the reference vector of several sheets]
    \label{thm:peps_sheet_stack_norm}
    \lean{TNLean.PEPS.EncodedFrame.one_sub_mul_le_norm_stackVec,
        TNLean.PEPS.EncodedFrame.one_sub_div_pow_le_norm_stackVec}
    \leanok
    \uses{def:peps_sheet_stack,thm:peps_frame_norm}
    If $\lVert\Omega\rVert=1$ and the $r$ total holes across $\mathcal F_1,\dots,\mathcal F_m$ each satisfy $\lVert(\id-P)\Omega\rVert\le\varepsilon$, then
    \begin{align}
        (1-r\varepsilon)\lVert\xi\rVert
         & \le\lVert\Omega_{\mathcal F_1}\otimes\cdots\otimes
        \Omega_{\mathcal F_m}\otimes\xi\rVert\le\lVert\xi\rVert. \notag
    \end{align}
    In particular, for $\varepsilon=L^{-60}$, $r\le CL^3$, and $\lVert\xi\rVert=1$, the norm lies in $[1-CL^{-57},1]$. Here $r\le CL^3$ is assumed as a hypothesis (in the source it bounds the frozen auxiliary holes; \texttt{06-geometry.tex}, lines 653--656 and 674--678).
\end{theorem}

\begin{proof}\leanok
    By Theorem~\ref{thm:peps_frame_norm}, frame $\mathcal F_i$ with $r_i$ holes satisfies $1-r_i\varepsilon\le\lVert\Omega_{\mathcal F_i}\rVert\le1$. Since $1-a\le x\le1$ and $(1-b)z\le y\le z$ imply $xy=y-(1-x)y\ge y-az\ge(1-a-b)z$, induction over $i=1,\dots,m$ yields the bound.
\end{proof}

\begin{definition}[The final regrouping]
    \label{def:peps_final_regrouping}
    \lean{TNLean.PEPS.EncodedFrame.trivialTag,
        TNLean.PEPS.EncodedFrame.physVec,
        TNLean.PEPS.EncodedFrame.freshMap,
        TNLean.PEPS.EncodedFrame.freshMap_apply,
        TNLean.PEPS.EncodedFrame.norm_freshMap_le,
        TNLean.PEPS.EncodedFrame.renameMap,
        TNLean.PEPS.EncodedFrame.renameMap_tmul,
        TNLean.PEPS.EncodedFrame.norm_renameMap_le,
        TNLean.PEPS.EncodedFrame.finOneOf,
        TNLean.PEPS.EncodedFrame.nonprivateCount_freshCircuit,
        TNLean.PEPS.EncodedFrame.nonprivateCount_siteCircuit,
        TNLean.PEPS.EncodedFrame.mem_nonsiteParties,
        TNLean.PEPS.EncodedFrame.nonsiteParties_dim,
        TNLean.PEPS.EncodedFrame.freshCircuit,
        TNLean.PEPS.EncodedFrame.siteCircuit,
        TNLean.PEPS.EncodedFrame.nonsiteParties,
        TNLean.PEPS.EncodedFrame.siteOrder,
        TNLean.PEPS.EncodedFrame.readoutPrivate,
        TNLean.PEPS.EncodedFrame.readoutCircuit}
    \leanok
    \uses{def:peps_sheet_stack}
    With the main frame $\mathcal M$ first in the sheet memory (Definition~\ref{def:peps_sheet_stack}), the final regrouping moves the raw registers of $\mathcal M$ ahead of its tag registers, privately renames each site's raw register at its owning party as that party's physical register, and privately prepares a $1$-dimensional physical register at each site-free party. Physical registers are ordered by sites then site-free parties; private registers consist of the tags of $\mathcal M$ followed by all other sheets and idle copies (\texttt{06-geometry.tex}, lines 48--51 and 669--681).
\end{definition}

\begin{theorem}[The final reference vector]
    \label{thm:peps_final_reference}
    \lean{TNLean.Lattice.PartyCircuit.OriginalCircuit.isEmpty_nonprivateLocations_of_nonprivateCount_eq_zero,
        TNLean.Lattice.PartyCircuit.OriginalCircuit.participants_subset_of_nonprivateCount_eq_zero,
        TNLean.Lattice.PartyCircuit.OriginalCircuit.isExpansionBounded_of_nonprivateCount_eq_zero,
        TNLean.Lattice.PartyCircuit.OriginalCircuit.isMonomialBounded_of_nonprivateCount_eq_zero,
        TNLean.PEPS.EncodedFrame.eq_trivialTag,
        TNLean.PEPS.EncodedFrame.rawProd_trivialTag,
        TNLean.PEPS.EncodedFrame.Frame.regRef_of_outer_eq_empty,
        TNLean.PEPS.EncodedFrame.familyPhysicalListBasis_eq_physVec,
        TNLean.PEPS.EncodedFrame.familyLabelledPhysicalBasis_eq_physVec,
        TNLean.PEPS.EncodedFrame.eval_freshCircuit,
        TNLean.PEPS.EncodedFrame.eval_siteCircuit,
        TNLean.PEPS.EncodedFrame.nodup_siteOrder,
        TNLean.PEPS.EncodedFrame.mem_siteOrder,
        TNLean.PEPS.EncodedFrame.nonprivateCount_readoutCircuit,
        TNLean.PEPS.EncodedFrame.eval_readoutCircuit_stackVec,
        TNLean.PEPS.EncodedFrame.readout_eval_readoutCircuit_stackVec,
        TNLean.PEPS.EncodedFrame.exists_readoutCircuit_norm,
        TNLean.PEPS.Approximation.DyadicProtocol.physicalConfigEquiv_apply_val,
        TNLean.PEPS.Approximation.DyadicProtocol.readout_eq_euclideanTensorVector}
    \leanok
    \uses{def:peps_final_regrouping, def:peps_dyadic_protocol}
    The final regrouping (Definition~\ref{def:peps_final_regrouping}) has $0$ nonprivate gates. If $\mathcal M$ has injective raw owner and empty physical outer sample at every hole, and the physical dimension is $q$ at site parties and $1$ elsewhere, then for any $\Omega$ and idle vector $\xi$ the regrouped reference vector $\Omega_{\mathcal M}\otimes\Omega_{\mathcal F_1}\otimes\cdots\otimes\Omega_{\mathcal F_m}\otimes\xi$ has physical-private coordinates $\Omega\otimes\Xi$ with
    \begin{align}
        \lVert\Xi\rVert & = \lVert\Omega_{\mathcal F_1}\otimes\cdots\otimes\Omega_{\mathcal F_m}\otimes\xi\rVert. \notag
    \end{align}
    If $\lVert\Omega\rVert=\lVert\xi\rVert=1$ and the $r\le CL^3$ holes of $\mathcal F_1,\dots,\mathcal F_m$ satisfy $\lVert(\id-P)\Omega\rVert\le L^{-60}$, then $1-CL^{-57}\le\lVert\Xi\rVert\le1$ and the protocol readout (Definition~\ref{def:peps_dyadic_protocol}) is $\Omega\otimes\Xi$ (\texttt{06-geometry.tex}, lines 43--51 and 669--681).
\end{theorem}

\begin{proof}\leanok
    \uses{thm:peps_sheet_stack_norm}
    Each empty-outer-sample hole has only the identity tag $\tau_0$ with identity raw operator, so $\Omega_{\mathcal M}=\sum_c\Omega(c)\,\lvert\tau_0\rangle\otimes\lvert c\rangle$. Moving $\lvert c\rangle$ to the front, renaming it to the site physical basis vector, and appending the $1$-dimensional nonsite basis vectors maps $\Omega_{\mathcal M}\otimes\Omega_{\mathcal F_1}\otimes\cdots\otimes\Omega_{\mathcal F_m}\otimes\xi$ by linearity to $\sum_c\Omega(c)\,\lvert c\rangle\otimes\Xi=\Omega\otimes\Xi$. Theorem~\ref{thm:peps_sheet_stack_norm} gives the norm bound on $\Xi$.
\end{proof}

\begin{theorem}[The initial step of the dyadic protocol]
    \label{thm:peps_dyadic_protocol_initial_step}
    \lean{TNLean.PEPS.Approximation.DyadicProtocol.exists_originalCircuit_initial,
        TNLean.PEPS.Approximation.rootSquare,
        TNLean.PEPS.Approximation.initialFrames,
        TNLean.PEPS.Approximation.length_initialFrames_le,
        TNLean.PEPS.Approximation.rawHolder,
        TNLean.PEPS.Approximation.mem_initialFrames,
        TNLean.PEPS.Approximation.DyadicProtocol.exists_originalCircuit_initialFrames}
    \leanok
    \uses{def:peps_sheet_stack, def:peps_dyadic_protocol,
        def:peps_dyadic_protocol_parties, def:peps_dyadic_protocol_schedule}
    Let $L\ge1$ and $\lVert\Omega\rVert=1$. The initial sheets consist of the hole-free main copy of $\Omega$ at the root-square party and one hole-free copy per auxiliary-initializing schedule change at its (arbitrary) future owner, totaling at most $1+4(\#\text{squares}+\#\text{junctions})$ sheets. If each party's initial private memory and vector are the memory and product reference vector of its held initial sheets, then one private map per party followed by a sheet reordering ($0$ nonprivate gates) maps the protocol's initial reference vector exactly to the reference vector of the initial sheets (\texttt{06-geometry.tex}, lines 21--23, 576, 649--652 and 658--667).
\end{theorem}

\begin{proof}\leanok
    \uses{thm:peps_sheet_raw_initialization, thm:peps_dyadic_protocol_schedule_counts}
    Every initial sheet is hole-free and held by a single party, so Theorem~\ref{thm:peps_sheet_raw_initialization} applies; Theorem~\ref{thm:peps_dyadic_protocol_schedule_counts} bounds the auxiliary initializations by $4(\#\text{squares}+\#\text{junctions})$.
\end{proof}

\begin{theorem}[The final step of the dyadic protocol]
    \label{thm:peps_dyadic_protocol_final_step}
    \lean{TNLean.PEPS.Approximation.DyadicProtocol.physicalDimension_eq_of_site_eq,
        TNLean.PEPS.Approximation.DyadicProtocol.physicalDimension_eq_one_of_site_eq,
        TNLean.PEPS.Approximation.DyadicProtocol.exists_reference_changes_of_readoutCircuit,
        TNLean.PEPS.Approximation.unitSquare,
        TNLean.PEPS.Approximation.toTriple_unitSquare,
        TNLean.PEPS.Approximation.unitSquare_injective,
        TNLean.PEPS.Approximation.isCellCenter_cellCenterPos,
        TNLean.PEPS.Approximation.blockIndex_one_cellCenterPos,
        TNLean.PEPS.Approximation.unitLevelFrame,
        TNLean.PEPS.Approximation.unitLevelFrame_owner,
        TNLean.PEPS.Approximation.unitLevelFrame_owner_injective,
        TNLean.PEPS.Approximation.DyadicProtocol.exists_reference_changes_unitLevel}
    \leanok
    \uses{def:peps_final_regrouping, def:peps_guide_frame,
        def:peps_dyadic_protocol_parties, def:peps_dyadic_protocol_schedule,
        def:peps_dyadic_protocol}
    Let $L\le 2^k<2L$ and $\lVert\Omega\rVert=1$. At the end of level $n=1$, each unit block of the root square carries its own label, so the site at $(i+\frac12,j+\frac12)$ is owned by the $(i,j)$ unit-square party injectively. Suppose the main frame has holes at grid corners of outer radius $<1/2$, the remaining sheets are one frozen auxiliary frame per auxiliary initialization (each with at most as many holes as true vertices of its final guide, using patches from the adaptive patch constraint at accuracy $L^{-60}$), and the idle registers carry a unit vector. Then the final regrouping (Definition~\ref{def:peps_final_regrouping}) produces protocol readout $\Omega\otimes\Xi$ with $1-214\,L^{-57}\le\lVert\Xi\rVert\le1$ (\texttt{06-geometry.tex}, lines 43--51, 653--656 and 669--681).
\end{theorem}

\begin{proof}\leanok
    \uses{thm:peps_final_reference, thm:peps_guide_frame_norm,
        lem:peps_guide_frame_patch_data, lem:peps_dyadic_frozen_hole_count,
        thm:peps_dyadic_repainting_valid_labels}
    Each site $(i+\frac12,j+\frac12)$ lies in unit block $(i,j)\subset[0,2^k]^2$, which has been repainted to the $(i,j)$ unit-square party's label (Theorem~\ref{thm:peps_dyadic_repainting_valid_labels}). By Theorem~\ref{thm:peps_guide_frame_norm}, all main-frame holes have empty physical outer sample; by Lemma~\ref{lem:peps_dyadic_frozen_hole_count}, the frozen auxiliary frames have at most $214\,L^3$ holes in total, each with projector error $\le L^{-60}$. Theorem~\ref{thm:peps_final_reference} yields the readout and norm bound.
\end{proof}

\begin{definition}[Bounded changes of the sheets]
    \label{def:peps_bounded_change}
    \lean{TNLean.PEPS.EncodedFrame.HasBoundedChange}
    \leanok
    \uses{def:peps_sheet_stack,def:lattice_bounded_change}
    An original circuit $C$ is a bounded change from $x$ to $y$ with participant set $S$, bounds $(r,K)$, and error $\varepsilon$ if $C$ has at most one nonprivate gate (with participants in $S$), every expansion monomial has $\le r$ pair effects, the monomial count and $\ell_1$ coefficient sum are $\le K$, and $\lVert Cx-y\rVert\le\varepsilon$. A bounded change of sheets from $\mathcal F_1,\dots,\mathcal F_m$ to $\mathcal F'_1,\dots,\mathcal F'_{m'}$ (with common idle registers $\ell$ and vector $\xi$) is a bounded change from $\Omega_{\mathcal F_1}\otimes\cdots\otimes\Omega_{\mathcal F_m}\otimes\xi$ to $\Omega_{\mathcal F'_1}\otimes\cdots\otimes\Omega_{\mathcal F'_{m'}}\otimes\xi$ (\cite[Proposition~7.1, items 2 and 4]{OpenAI2026PolynomialPEPS}; \texttt{06-geometry.tex}, lines 25--41 and 572--583).
\end{definition}

\begin{theorem}[Operations that preserve bounded changes]
    \label{thm:peps_bounded_change_closure}
    \lean{TNLean.PEPS.EncodedFrame.pairRef_eq_twoLayoutIso_symm,
        TNLean.PEPS.EncodedFrame.HasBoundedChange.mono,
        TNLean.PEPS.EncodedFrame.hasBoundedChange_of_perm,
        TNLean.PEPS.EncodedFrame.hasBoundedChange_refl,
        TNLean.PEPS.EncodedFrame.HasBoundedChange.perm_left,
        TNLean.PEPS.EncodedFrame.HasBoundedChange.perm_right,
        TNLean.PEPS.EncodedFrame.isBoundedChange_liftOne,
        TNLean.PEPS.EncodedFrame.isBoundedChange_liftTwo,
        TNLean.PEPS.EncodedFrame.hasBoundedChange_liftOne,
        TNLean.PEPS.EncodedFrame.hasBoundedChange_liftTwo,
        TNLean.PEPS.EncodedFrame.hasBoundedChange_activate,
        TNLean.PEPS.EncodedFrame.hasBoundedChange_perm_liftOne,
        TNLean.PEPS.EncodedFrame.exists_isBoundedChange_of_frame_eq}
    \leanok
    \uses{def:peps_bounded_change,thm:peps_sheet_reordering,thm:peps_identity_change,
        thm:lattice_bounded_change_closure}
    Let $\lVert\Omega\rVert\le1$ and $\lVert\xi\rVert\le1$.
    \begin{enumerate}
        \item Sheet permutations and equal-frame transitions are bounded changes for any $S,r,K$ and $\varepsilon\ge0$, and composing a bounded change with a sheet permutation on either side preserves $(S,r,K,\varepsilon)$.
        \item A bounded change $\Omega_{\mathcal F}\mapsto\Omega_{\mathcal F'}$ (or $\Omega_{\mathcal F_1}\otimes\Omega_{\mathcal F_2}\mapsto\Omega_{\mathcal F'_1}\otimes\Omega_{\mathcal F'_2}$) lifts to a bounded change of a sheet stack replacing $\mathcal F$ by $\mathcal F'$ (or adjacent $\mathcal F_1,\mathcal F_2$ by $\mathcal F'_1,\mathcal F'_2$) with the same $(S,r,K,\varepsilon)$.
        \item An allowed word on $S$ with $\le1$ pair source is a bounded change with its own error for any $r\ge0,K\ge1$; an allowed chain on $S$ with $0$ pair sources and $\le1$ pair effect is a bounded change for any $r,K\ge1$.
    \end{enumerate}
\end{theorem}

\begin{proof}\leanok
    \uses{thm:peps_sheet_reordering,
        thm:peps_original_produced_resource_bounds}
    Sheet permutations are exact exchange words with $0$ nonprivate gates (Theorem~\ref{thm:peps_sheet_reordering}), and original circuits are contractions. Spectator frames have reference vectors of norm $\le\lVert\Omega\rVert\le1$. Allowed words and chains convert to original circuits with one two-party single-monomial gate per pair source or pair effect (Theorem~\ref{thm:peps_original_produced_resource_bounds}).
\end{proof}

\begin{theorem}[Frames under maps agreeing off the placeholder]
    \label{thm:peps_guide_frame_owner_agreement}
    \lean{TNLean.PEPS.Approximation.guideFrame_eq_of_owner_eq}
    \leanok
    \uses{def:peps_guide_frame, def:peps_dyadic_protocol_parties}
    Let $g$ be a guide taking only valid labels and avoiding the placeholder at every site, and let party maps $\rho,\rho'$ satisfy $\rho(\ell)=\rho'(\ell)$ for every valid non-placeholder label $\ell$. For any common label rank, finite set $V$ of true vertices of $g$, and square patches on $V$ with pairwise disjoint physical outer samples, the guide frames of $g$ under $\rho$ and $\rho'$ coincide.
\end{theorem}

\begin{proof}\leanok
    \uses{def:peps_guide_frame, lem:peps_guide_frame_tags, lem:peps_dyadic_tags}
    Site labels are non-placeholder, so raw owners agree; canonical tag holders on $V$ are also valid non-placeholder labels by Lemmas~\ref{lem:peps_guide_frame_tags} and~\ref{lem:peps_dyadic_tags}, so tag owners and hole lists agree as well.
\end{proof}

\begin{theorem}[Switching the party map between repaintings]
    \label{thm:peps_protocol_party_map_switch}
    \lean{TNLean.PEPS.Approximation.labelParty_eq_of_ne_guidePlaceholder,
        TNLean.PEPS.Approximation.guideFrame_labelParty_eq,
        TNLean.PEPS.Approximation.exists_isBoundedChange_labelParty_switch,
        TNLean.PEPS.Approximation.openRoot,
        TNLean.PEPS.Approximation.blockGuide_protocolLevelLabels_ne_guidePlaceholder,
        TNLean.PEPS.Approximation.IsDisplacedReading.ne_guidePlaceholder,
        TNLean.PEPS.Approximation.cellCenterPos_mem_openRoot,
        TNLean.PEPS.Approximation.guideFrame_levelLabels_labelParty_eq,
        TNLean.PEPS.Approximation.exists_isBoundedChange_levelLabels_labelParty_switch}
    \leanok
    \uses{def:peps_dyadic_protocol_parties}
    For any guide taking only the placeholder and hierarchy block labels, with no placeholder at any site and all holes centered at true vertices, the encoded frame is identical under the party maps of any two repaintings, so switching party maps is the identity bounded change for any $(S,r,K,\varepsilon)$. At the start and end of each repainting at level $2^{j'}$ ($j'<k$), every reading of the level block guide takes a block label on $(0,2^k)^2$ (which contains all sites $(i+\frac12,j+\frac12)$ for $0\le i,j<L\le2^k$), so the party map can be switched and the fresh placeholder party owns no site or tag before or after the repainting (\texttt{06-geometry.tex}, lines 82--100, 119--125 and 627--630).
\end{theorem}

\begin{proof}\leanok
    \uses{thm:peps_guide_frame_owner_agreement, thm:peps_dyadic_guide_owners,
        lem:peps_dyadic_sampled_guides, thm:peps_bounded_change_closure}
    Every hierarchy block label is mapped to its block party by every repainting's party map, and tag holders at true vertices are non-placeholder incident labels, so Theorem~\ref{thm:peps_guide_frame_owner_agreement} and Theorem~\ref{thm:peps_bounded_change_closure} apply. Every point of $(0,2^k)^2$ has a neighborhood in $(0,2^k)^2$ where the level block guide takes block labels, which displaced readings preserve (Lemma~\ref{lem:peps_dyadic_sampled_guides}).
\end{proof}

\begin{theorem}[The operations of the protocol as bounded changes]
    \label{thm:peps_protocol_bounded_changes}
    \lean{TNLean.PEPS.Approximation.DyadicSquare.index,
        TNLean.PEPS.Approximation.DyadicSquare.inRoot_index,
        TNLean.PEPS.Approximation.eventUses_labelParty_protocolLevelLabels,
        TNLean.PEPS.Approximation.eventUses_labelParty_newParty,
        TNLean.PEPS.Approximation.eventUses_labelParty_guidePlaceholder,
        TNLean.PEPS.Approximation.abs_sub_le_one_of_mem_blockWindow,
        TNLean.PEPS.Approximation.coe_subset_eventUses_of_window,
        TNLean.PEPS.Approximation.coe_blockParties_subset_eventUses,
        TNLean.PEPS.Approximation.IsProtocolConstants.lt_one_div_four_thousand,
        TNLean.PEPS.Approximation.isBoundedChange_of_step,
        TNLean.PEPS.Approximation.directRepainting_isBoundedChange,
        TNLean.PEPS.Approximation.isBoundedChange_of_birth,
        TNLean.PEPS.Approximation.isBoundedChange_of_death,
        TNLean.PEPS.Approximation.isBoundedChange_of_exchange}
    \leanok
    \uses{thm:peps_bounded_change_closure,
        def:peps_dyadic_protocol_parties,
        thm:peps_dyadic_direct_repainting_constants,
        def:peps_dyadic_block_guide,
        def:peps_dyadic_hierarchy}
    Every hierarchy square has its index in the root square of its level, and the protocol constants satisfy $\epsilon_0<1/4000$. Under the party map of a repainting of $R$, the labels of the $3\times3$ window around $R$ at any stage of its level pass, the new label of $R$, and the placeholder all map to parties used by the repainting of $R$. Consequently, under Theorem~\ref{thm:peps_dyadic_direct_repainting_constants}, direct repainting of $R$ is a bounded change with participants used by $R$, error $(2b+1)L^{-20}$, and $\le 2b$ pair sources per monomial ($b$ the protocol patch bound). Likewise, each point treatment, treatment removal, resize, bulk birth, bulk death, or lens exchange (with $0$ nonprivate gates) is a bounded change with participants among the parties of its step theorem (\texttt{06-geometry.tex}, lines 91--100, 328--541 and 559--596).
\end{theorem}

\begin{proof}\leanok
    \uses{thm:peps_dyadic_direct_repainting_constants,
        thm:peps_dyadic_protocol_parties,
        thm:peps_dyadic_repainting_valid_labels,
        thm:peps_dyadic_active_label_parties,
        def:peps_dyadic_used_parties,
        def:peps_dyadic_protocol_parties,
        thm:peps_dyadic_block_ratio,
        thm:peps_guide_frame_resize_treatment_gates,
        lem:peps_dyadic_active_labels,
        thm:peps_bounded_change_closure,
        thm:peps_protocol_constants,
        def:peps_bounded_change,
        def:peps_sheet_stack,
        def:peps_dyadic_block_guide}
    The root-index bound follows from Definition~\ref{def:peps_dyadic_hierarchy}, and $2\epsilon_0(a_0+1)\le a_0/2$ with $a_0=1/5000$ gives $\epsilon_0\le 1/20004<1/4000$ (Theorem~\ref{thm:peps_protocol_constants}). By Lemma~\ref{lem:peps_dyadic_active_labels} and Theorem~\ref{thm:peps_dyadic_active_label_parties}, labels in the $3\times3$ window around $R$ map to parties used by $R$ (Definition~\ref{def:peps_dyadic_used_parties}). Theorem~\ref{thm:peps_bounded_change_closure} and Theorem~\ref{thm:peps_dyadic_direct_repainting_constants} then convert each step to a bounded change while preserving the $2b$ pair-source bound.
\end{proof}

\begin{lemma}[The parties of a junction resize]
    \label{lem:peps_protocol_junction_resize_parties}
    \lean{TNLean.PEPS.Approximation.DyadicJunction.index,
        TNLean.PEPS.Approximation.eventUses_labelParty_protocolLevelLabels_junction,
        TNLean.PEPS.Approximation.coe_subset_eventUses_of_junctionWindow,
        TNLean.PEPS.Approximation.coe_subset_eventUses_of_junctionWindow_placeholder}
    \leanok
    \uses{def:peps_dyadic_protocol_parties, def:peps_dyadic_used_parties}
    At the end of the level-$2^j$ pass, consider the resize at grid corner $(2^ju,2^jv)$ under the party map of any repainting $R$. Every root-square block in $(u-1,v-1)+\{-1,0,1\}^2$ is held by a block party used by the resize, while any block outside the root square carries the placeholder party of $R$ (\texttt{06-geometry.tex}, lines 91--100, 127--131, 529--541 and 594).
\end{lemma}

\begin{proof}\leanok
    Each root-square block in the window carries its scale-$j$ block party with indices in $[u-2,u]\times[v-2,v]$, which belongs to the resize's used parties by Definition~\ref{def:peps_dyadic_used_parties}.
\end{proof}

\begin{theorem}[Bounded changes along the schedule]
    \label{thm:peps_protocol_bounded_trajectory}
    \lean{TNLean.PEPS.EncodedFrame.exists_steps_of_hasBoundedChange,
        TNLean.PEPS.Approximation.norm_trajectory_sub_stackVec_le,
        TNLean.PEPS.Approximation.card_participation_le_of_eventUses}
    \leanok
    \uses{thm:peps_bounded_change_closure, def:peps_dyadic_protocol_schedule,
        thm:peps_dyadic_protocol_lifetime}
    Given a sequence of frame stacks joined by bounded changes of error $\le CL^{-20}$ along a schedule of length $\le 5040L^3$ (for $1\le L$ and $2^k<2L$), the trajectory started at the initial reference vector stays within $5040\,CL^{-17}$ of the reference vector at every stage. If each schedule change has $\le1$ nonprivate gate with participants used by its event, every party participates in at most $64\cdot189$ nonprivate gates (\texttt{06-geometry.tex}, lines 25--41 and 572--631).
\end{theorem}

\begin{proof}\leanok
    \uses{thm:peps_dyadic_protocol_lifetime}
    Contractivity of original circuits sums the per-step errors over the $\le5040L^3$ changes, and $\le1$ gate per change gives $\le189$ gates per event, so Theorem~\ref{thm:peps_dyadic_protocol_lifetime} bounds the per-party gate participation by $64\cdot189$.
\end{proof}

\begin{theorem}[A protocol from a bounded trajectory of sheets]
    \label{thm:peps_protocol_from_bounded_trajectory}
    \lean{TNLean.Lattice.PartyCircuit.OriginalCircuit.chainLayout,
        TNLean.Lattice.PartyCircuit.OriginalCircuit.chainStep,
        TNLean.Lattice.PartyCircuit.OriginalCircuit.chainVec,
        TNLean.Lattice.PartyCircuit.OriginalCircuit.chainStep_induction,
        TNLean.Lattice.PartyCircuit.OriginalCircuit.norm_chainStep_eval_sub_le,
        TNLean.PEPS.Approximation.card_eventUses_le,
        TNLean.PEPS.Approximation.pow_le_paddedSide,
        TNLean.PEPS.Approximation.finiteDimensional_of_mem_stackRegs,
        TNLean.PEPS.Approximation.exists_isDistributionProtocol_of_frames}
    \leanok
    \uses{def:peps_dyadic_protocol, thm:peps_dyadic_protocol_initial_step,
        thm:peps_dyadic_protocol_final_step, thm:peps_protocol_bounded_trajectory,
        thm:peps_dyadic_protocol_locality, thm:peps_dyadic_protocol_parties}
    Every hierarchy event uses at most $19$ parties. Let $L\ge2$, $L\le2^k<2L$, $\lVert\Omega\rVert=1$, and suppose frame stacks on the hierarchy parties of $[0,2^k]^2$ are given at every stage such that:
    \begin{enumerate}
        \item The initial frames have no holes, and each is held entirely by one
              party.
        \item Each change of the schedule is a bounded change from the reference
              vector of its stage to that of the next stage, with participants
              used by the event of the change, at most $r$ pair effects and at most
              $r$ pair sources in each monomial, at most $K\le C_KL^{k'}$ monomials
              and coefficient sum, and error at most $C_EL^{-20}$.
        \item At the end of the schedule the main frame is read by the unit-square
              parties and has trivial holes, every hole of the other sheets
              satisfies the frame constraint at accuracy $L^{-60}$, and the other
              sheets carry at most $C_EL^3$ holes.
    \end{enumerate}
    If $C_N\ge5042$, the private initialization, the changes of the schedule and the
    final regrouping form a protocol with the assembly bounds of
    Definition~\ref{def:peps_dyadic_protocol} for the count constant $C_N$, the
    expansion constant $C_K$, the error constant $C_E$, and the constants
    $b=64\cdot189$, $k'$ and $r$, including the bound $r$ on the pair sources
    of each monomial (\texttt{06-geometry.tex}, lines 12--53 and 572--681). This is not
    Proposition~\ref{prop:peps_dyadic_protocol}: the frames along the schedule are
    assumed here.
\end{theorem}

\begin{proof}\leanok
    \uses{thm:peps_sheet_raw_initialization,
        thm:peps_dyadic_protocol_lifetime,
        thm:peps_dyadic_protocol_schedule_counts,
        thm:peps_original_produced_resource_bounds}
    A square repainting uses $\le 18$ neighbor block parties plus its placeholder ($19$ total); a junction resize uses $\le 9$. Combining exact private initialization (Theorem~\ref{thm:peps_dyadic_protocol_initial_step}), the $\le 5040L^3$ schedule changes (Theorems~\ref{thm:peps_protocol_bounded_trajectory}, \ref{thm:peps_dyadic_protocol_schedule_counts}, and~\ref{thm:peps_dyadic_protocol_locality}), and the exact final regrouping (Theorem~\ref{thm:peps_dyadic_protocol_final_step}) yields all bounds of Definition~\ref{def:peps_dyadic_protocol}.
\end{proof}

\begin{theorem}[Joining the changes of the events]
    \label{thm:peps_protocol_event_segments}
    \lean{TNLean.PEPS.EncodedFrame.HasSourceBoundedChange,
        TNLean.PEPS.EncodedFrame.HasSourceBoundedChange.hasBoundedChange,
        TNLean.PEPS.EncodedFrame.HasSourceBoundedChange.mono,
        TNLean.PEPS.EncodedFrame.hasSourceBoundedChange_of_perm,
        TNLean.PEPS.EncodedFrame.hasSourceBoundedChange_refl,
        TNLean.PEPS.EncodedFrame.HasSourceBoundedChange.perm_left,
        TNLean.PEPS.EncodedFrame.HasSourceBoundedChange.perm_right,
        TNLean.PEPS.EncodedFrame.isSourceBounded_liftOne,
        TNLean.PEPS.EncodedFrame.isSourceBounded_liftTwo,
        TNLean.PEPS.EncodedFrame.hasSourceBoundedChange_liftOne,
        TNLean.PEPS.EncodedFrame.hasSourceBoundedChange_liftTwo,
        TNLean.PEPS.EncodedFrame.hasSourceBoundedChange_activate,
        TNLean.PEPS.EncodedFrame.hasSourceBoundedChange_perm_liftOne,
        TNLean.PEPS.Approximation.exists_frames_of_segments,
        TNLean.PEPS.Approximation.exists_isDistributionProtocol_of_hasSourceBoundedChange}
    \leanok
    \uses{def:peps_bounded_change, def:peps_circuit_source_bound,
        thm:peps_protocol_from_bounded_trajectory}
    A source-bounded change (a bounded change with $\le r_s$ pair sources per monomial, Definition~\ref{def:peps_circuit_source_bound}) is preserved by enlarging $(S,r,K,r_s,\varepsilon)$, by composing with sheet permutations or identity changes, and (when $\lVert\Omega\rVert,\lVert\xi\rVert\le1$) by lifting from one or two adjacent sheets to a larger stack. Consequently, if each event in the schedule carries a segment of source-bounded changes connecting its starting stack to the next event's starting stack with participants used by that event, these segments concatenate into a full schedule trajectory, and under the remaining hypotheses of Theorem~\ref{thm:peps_protocol_from_bounded_trajectory} its conclusion holds (\texttt{06-geometry.tex}, lines 82--100, 127--131 and 572--681; \texttt{04-compression.tex}, lines 143--147).
\end{theorem}

\begin{proof}\leanok
    \uses{thm:peps_bounded_change_closure, thm:peps_sheet_reordering,
        lem:peps_circuit_source_bound_monomials, thm:peps_protocol_from_bounded_trajectory}
    By Theorem~\ref{thm:peps_bounded_change_closure} and Lemma~\ref{lem:peps_circuit_source_bound_monomials}, tensoring with untouched spectator registers adds no pair sources and does not increase the error when $\lVert\Omega\rVert,\lVert\xi\rVert\le1$. Inducting along the event list concatenates the segments, and Theorem~\ref{thm:peps_protocol_from_bounded_trajectory} concludes.
\end{proof}

\begin{definition}[The main frame between repaintings]
    \label{def:peps_protocol_level_frame}
    \lean{TNLean.PEPS.Approximation.sampleGuide,
        TNLean.PEPS.Approximation.levelGuide,
        TNLean.PEPS.Approximation.levelVertices,
        TNLean.PEPS.Approximation.levelFrame}
    \leanok
    \uses{def:peps_guide_frame, def:peps_dyadic_protocol_parties,
        def:peps_adaptive_patch_constraint, lem:peps_level_true_vertices,
        def:peps_dyadic_block_guide, thm:peps_dyadic_repainting_valid_labels,
        thm:peps_guide_frame_supplied_gates, def:peps_frame}
    For a plane labelling $G$, its sampled guide at $p$ is $\lim_{\varepsilon\to0^+}G(p+\varepsilon(1,1))$ when eventually constant, and $G(p)$ otherwise. Given hierarchy depth $k$, level $0\le j\le k$, repainted blocks $\mathcal D$, a label rank, a reference repainting party map, an adaptive patch supply $P$ with parameters $(B_0,C,c,60,L)$, and a radius function $0<\rho(z)\le B_0\log L$, the level guide is the sampled level-$j$ block guide; its true vertices form a finite set of grid corners in $[0,2^k]^2$. The level frame is the guide frame (Definition~\ref{def:peps_guide_frame}) with supplied patches of radius $\rho(z)$ at those true vertices if their outer physical samples are pairwise disjoint, and the root-owned raw frame otherwise (\texttt{06-geometry.tex}, lines 69--125).
\end{definition}

\begin{theorem}[Properties of the main frame between repaintings]
    \label{thm:peps_protocol_level_frame_facts}
    \lean{TNLean.PEPS.Approximation.IsDisplacedReading.eq_sampleGuide,
        TNLean.PEPS.Approximation.IsRayConstant.isDisplacedReading_sampleGuide,
        TNLean.PEPS.EncodedFrame.hasSourceBoundedChange_main,
        TNLean.Lattice.perm_map_update_cons,
        TNLean.Lattice.perm_map_cons,
        TNLean.PEPS.EncodedFrame.hasSourceBoundedChange_aux,
        TNLean.PEPS.EncodedFrame.hasSourceBoundedChange_mainAux,
        TNLean.PEPS.Approximation.isDisplacedReading_levelGuide,
        TNLean.PEPS.Approximation.mem_levelVertices,
        TNLean.PEPS.Approximation.levelFrame_eq,
        TNLean.PEPS.Approximation.guideFrame_congr_guide,
        TNLean.PEPS.Approximation.levelFrame_eq_of_levelGuide_eq,
        TNLean.PEPS.Approximation.levelFrame_congr,
        TNLean.PEPS.Approximation.pairwise_disjoint_levelHoles,
        TNLean.PEPS.Approximation.isTrueVertex_levelGuide_insert_iff}
    \leanok
    \uses{def:peps_protocol_level_frame, def:peps_sheet_stack,
        thm:peps_protocol_event_segments, def:peps_circuit_source_bound,
        lem:peps_dyadic_sampled_guides, lem:peps_dyadic_readings_exist}
    \begin{enumerate}
        \item Every $(1,1)$-displaced reading of $G$ equals its sampled guide, and when $G$ is ray-constant (in particular for block guides) the sampled guide is a $(1,1)$-reading.
        \item When $q\ge1$ and $\lVert\Omega\rVert,\lVert\xi\rVert\le1$, a source-bounded change of the main frame, of one auxiliary frame (indexed in a duplicate-free list), or of both together lifts to the full sheet memory with the same bounds $(S,r,K,r_s,\varepsilon)$ via sheet permutations moving the active auxiliary frame adjacent to the main frame.
        \item The finite vertex set of Definition~\ref{def:peps_protocol_level_frame} is the exact set of true vertices of the level guide; equal sampled guides with radius functions agreeing on all true vertices produce equal level frames; if $\rho(z)\le r$ with $4r<2^j$, the holes have pairwise disjoint outer physical samples; and repainting one block changes true-vertex status only at its four corners (\texttt{06-geometry.tex}, lines 69--125 and 507--511).
    \end{enumerate}
\end{theorem}

\begin{proof}\leanok
    \uses{def:peps_protocol_level_frame, def:peps_guide_frame,
        thm:peps_guide_frame_patch_congruence,
        lem:peps_dyadic_readings_exist, thm:peps_protocol_event_segments,
        lem:peps_level_true_vertices, lem:peps_dyadic_choice_K0,
        thm:peps_guide_frame_norm, thm:peps_guide_frame_supplied_gates,
        def:peps_dyadic_block_guide, lem:peps_dyadic_direct_repainting,
        lem:peps_dyadic_incident_change, lem:peps_dyadic_sampled_guides,
        thm:peps_dyadic_temporary_holes, lem:peps_dyadic_active_labels,
        def:peps_dyadic_incident_labels, thm:peps_dyadic_repainting_valid_labels}
    Ray-constancy of block guides (Lemma~\ref{lem:peps_dyadic_readings_exist}) gives (1), and Theorem~\ref{thm:peps_protocol_event_segments} gives (2). For (3), Lemma~\ref{lem:peps_level_true_vertices} and Theorem~\ref{thm:peps_guide_frame_patch_congruence} identify the true vertices and frames; grid corners at level $j$ have spacing $\ge 2^j>4r$, so outer samples of radius $2r$ are disjoint; and outside the four corners of a repainted block the block guide is locally unchanged or interior to the block, preserving true-vertex status.
\end{proof}

\begin{theorem}[Direct repaintings and resizes as changes of the protocol]
    \label{thm:peps_protocol_level_events}
    \lean{TNLean.PEPS.Approximation.rewriteGateBound,
        TNLean.PEPS.Approximation.IsProtocolConstants.holeRadius_bounds,
        TNLean.PEPS.Approximation.IsProtocolConstants.resizeRadius_bounds,
        TNLean.PEPS.Approximation.directRepainting_hasSourceBoundedChange,
        TNLean.PEPS.Approximation.levelVertices_self,
        TNLean.PEPS.Approximation.resizeRadius,
        TNLean.PEPS.Approximation.DyadicJunction.rootNbr,
        TNLean.PEPS.Approximation.DyadicJunction.rootNbr_mem,
        TNLean.PEPS.Approximation.DyadicJunction.inRoot_rootNbr,
        TNLean.PEPS.Approximation.resizeParty,
        TNLean.PEPS.Approximation.eventUses_resizeParty,
        TNLean.PEPS.Approximation.resize_hasSourceBoundedChange_of_le,
        TNLean.PEPS.Approximation.exists_resize_hasSourceBoundedChange}
    \leanok
    \uses{def:peps_protocol_level_frame, thm:peps_protocol_constants,
        def:peps_dyadic_protocol_parties, thm:peps_protocol_event_segments,
        def:peps_dyadic_hierarchy, def:peps_adaptive_patch_constraint,
        thm:peps_dyadic_direct_repainting_constants, def:peps_circuit_source_bound,
        def:peps_sheet_stack}
    Let $2\le L\le 2^k$, $q\ge1$, $\lVert\Omega\rVert=1$, protocol constants $(K_0,\epsilon_0,\nu,a,D,B_0)$ as in Theorem~\ref{thm:peps_protocol_constants} with patch bound $b$, and let $P$ be an adaptive patch supply for $\Omega$ with parameters $(B_0,C,c,60,L)$. Fix a label rank, reference party map, spectator frames, and integers $C'\ge\max(C,1)$, $c'\ge c$. Then $t=D\log L$ and $h_n=\epsilon_0\min(n,t)$ satisfy $0<h_n\le B_0\log L$ for all $n>0$, and with
    \begin{align}
        K_{\mathrm{rw}} & =(C'+1)^{3b}
        \bigl(16b^2(C'+1)^{6b}+1\bigr)^{2b}
        L^{3bc'+4b(3bc'+30)},
        \label{eq:peps_protocol_level_rewrite_bound}
    \end{align}

    For all such data, assume also the direct-repainting ratio condition of
    Theorem~\ref{thm:peps_dyadic_direct_repainting_constants}. Repainting a square of
    the hierarchy at level $j<k$, of side $n=2^j\le K_0t$, from any set of already
    repainted blocks to that set with this square inserted, with all holes of radius
    $h_n$, gives a change from its old level frame to its new level frame with the same
    spectators. The change uses only participants used by the event, has at most $2b$
    pair effects and $2b$ pair sources in each monomial, at most $K_{\mathrm{rw}}$
    monomials and coefficient sum, and error $(2b+1)L^{-20}$, with empty external
    spectator layout and scalar $1$.

    For $0<\epsilon_0\le1/16$, let $\nu'=\min(\epsilon_0,1)/8$. For every
    $0<\nu\le\nu'$, every hierarchy $k$ and every choice of the remaining common data, the following
    resize has the same bounds. Let $J$ be a junction of the hierarchy of scale
    $1\le j\le k$, put $n=2^{j-1}$, and take any set $A\subseteq\mathbb R^2$. At the
    end of the pass to level $2^j$, give the holes radius $h_n$ at points of $A$ and
    radius $h_{2n}$ elsewhere. Replace $A$ by $A\cup\{z_J\}$, where $z_J$ is the corner
    of $J$. This change has the stated bounds with the same spectators and uses only
    participants used by the resize. No direct-repainting ratio condition is required
    for this resize assertion. If the two radii agree, or the corner is not a true
    vertex, the frame is unchanged. The root-level guide has no true vertex.
    If $Q_J$ is the junction index, its chosen offset $\delta_J$ belongs
    to $\{-1,0,1\}^2$, and $\delta_J+Q_J-(1,1)$ indexes a block inside the root
    square. Assigning the placeholder to that block's party makes every party of
    the translated nine-block window a participant used by the resize
    (\texttt{06-geometry.tex}, lines 91--133, 491, 507--541, 564--567 and 572--596).
\end{theorem}

\begin{proof}\leanok
    \uses{thm:peps_protocol_party_map_switch, thm:peps_guide_frame_owner_agreement,
        thm:peps_protocol_level_frame_facts, thm:peps_protocol_bounded_changes,
        thm:peps_guide_frame_supplied_gates, lem:peps_protocol_junction_resize_parties,
        thm:peps_protocol_event_segments, thm:peps_protocol_constants,
        thm:peps_dyadic_guide_owners, def:peps_dyadic_chamber_sampling,
        lem:peps_dyadic_incident_change, def:peps_dyadic_block_guide,
        thm:peps_dyadic_repainting_valid_labels, thm:peps_dyadic_protocol_parties,
        lem:peps_dyadic_floor, thm:peps_dyadic_block_ratio}
    Since $0<\epsilon_0<1\le K_0$ and $K_0D\le B_0$, we have $0<h_n\le\epsilon_0D\log L\le B_0\log L$, and $4\epsilon_0<1$ ensures disjoint outer physical samples at spacing $s$ (Theorem~\ref{thm:peps_protocol_level_frame_facts}). Switching to the event's party map via Theorem~\ref{thm:peps_protocol_party_map_switch} (and Lemma~\ref{lem:peps_protocol_junction_resize_parties} for the junction placeholder), the supplied-gate bounds of Theorems~\ref{thm:peps_guide_frame_supplied_gates} and~\ref{thm:peps_protocol_bounded_changes} lifted by Theorem~\ref{thm:peps_protocol_level_frame_facts} yield both claims.
\end{proof}

\begin{theorem}[The main frame along the schedule]
    \label{thm:peps_protocol_main_frame_schedule}
    \lean{TNLean.PEPS.Approximation.ProtocolEvent.two_le_rank,
        TNLean.PEPS.Approximation.ProtocolEvent.rank_le,
        TNLean.PEPS.Approximation.exists_isScheduled_rank,
        TNLean.PEPS.Approximation.isScheduled_getElem,
        TNLean.PEPS.Approximation.take_succ_eventSchedule,
        TNLean.PEPS.Approximation.mem_take_of_rank_lt,
        TNLean.PEPS.Approximation.rank_le_of_mem_take,
        TNLean.PEPS.Approximation.rank_getElem_succ_le,
        TNLean.PEPS.Approximation.exists_index_of_isTrueVertex_levelGuide,
        TNLean.PEPS.Approximation.resizedCorners,
        TNLean.PEPS.Approximation.repaintedBlocks,
        TNLean.PEPS.Approximation.levelGuide_eq_univ,
        TNLean.PEPS.Approximation.levelGuide_succ_univ,
        TNLean.PEPS.Approximation.eventMain,
        TNLean.PEPS.Approximation.finalMain,
        TNLean.PEPS.Approximation.mainState,
        TNLean.PEPS.Approximation.mainState_of_lt,
        TNLean.PEPS.Approximation.mainState_of_le,
        TNLean.PEPS.Approximation.levelGuide_eq_of_mem_open,
        TNLean.PEPS.Approximation.eventMain_inr_succ,
        TNLean.PEPS.Approximation.eventMain_inl_succ,
        TNLean.PEPS.Approximation.mainState_zero,
        TNLean.PEPS.Approximation.finalMain_owner,
        TNLean.PEPS.Approximation.finalMain_holes,
        TNLean.PEPS.Approximation.ProtocolEvent.rank_inl,
        TNLean.PEPS.Approximation.ProtocolEvent.rank_inr,
        TNLean.PEPS.Approximation.rank_getElem_le_of_le,
        TNLean.PEPS.Approximation.resizedCorners_succ_inr,
        TNLean.PEPS.Approximation.repaintedBlocks_succ_inl,
        TNLean.PEPS.Approximation.holeRadius_mem,
        TNLean.PEPS.Approximation.resizeRadius_mem,
        TNLean.PEPS.Approximation.squareOfIndex,
        TNLean.PEPS.Approximation.squareOfIndex_index,
        TNLean.PEPS.Approximation.junctionOfIndex,
        TNLean.PEPS.Approximation.junctionOfIndex_index,
        TNLean.PEPS.Approximation.getElem_notMem_take,
        TNLean.PEPS.Approximation.notMem_take_succ_of_rank,
        TNLean.PEPS.Approximation.rank_getElem_last,
        TNLean.PEPS.Approximation.rank_getElem_zero,
        TNLean.PEPS.Approximation.length_eventSchedule_pos,
        TNLean.PEPS.Approximation.inRoot_subset_repaintedBlocks,
        TNLean.PEPS.Approximation.frame_eq_of_owner_holes,
        TNLean.PEPS.Approximation.finalMain_pairwise}
    \leanok
    \uses{def:peps_protocol_level_frame, def:peps_dyadic_protocol_schedule,
        def:peps_adaptive_patch_constraint, def:peps_frame,
        def:peps_dyadic_hierarchy, lem:peps_dyadic_floor}
    In a depth-$k$ hierarchy ($0\le j\le k$), a level-$j$ square repainting has rank $2(k-j)+1$ and a scale-$j$ junction resize has rank $2(k-\max\{j-1,0\})$. The event schedule is duplicate-free with nondecreasing ranks covering $[2,2k+1]$ (when $k\ge1$). After a scale-$j$ resize or level-$j$ repainting, its corner or block index is added to the resized corners $\mathcal C_j$ or repainted blocks $\mathcal D_j$; when the next event has strictly larger rank (or at schedule end), all root-square blocks at level $j$ lie in $\mathcal D_j$.

    Let $k\ge1$, $L\le 2^k$, $q\ge1$, $\epsilon_0,t>0$ with $\epsilon_0t\le B_0\log L$, and $h_n=\epsilon_0\min(n,t)\in(0,B_0\log L]$. Fix an adaptive patch supply for $\Omega$ with parameters $(B_0,C,c,60,L)$, a label rank, and a reference party map. Define the main frame (Definition~\ref{def:peps_protocol_level_frame}) before each event as:
    \begin{enumerate}
        \item before a scale-$j$ resize: the level-$j$ frame on uniform $2^j$-blocks with hole radius $h_{2^{j-1}}$ on $\mathcal C_j$ and $h_{2^j}$ elsewhere;
        \item before a level-$j$ repainting: the level-$j$ frame at $\mathcal D_j$ with all hole radii $h_{2^j}$;
        \item at or after schedule end: the unit-square level frame.
    \end{enumerate}
    Then each resize or repainting transitions the main frame to that before the next event by inserting its corner or square; the initial main frame at index $0$ is the root-owned raw copy; and when $4\epsilon_0<1$, the final main frame has each site owned by its unit-square party and empty physical outer squares at all holes (\texttt{06-geometry.tex}, lines 21--23, 82--133 and 658--673).
\end{theorem}

\begin{proof}\leanok
    \uses{thm:peps_dyadic_protocol_schedule_counts,
        lem:peps_dyadic_direct_repainting, lem:peps_dyadic_unmodified_vertices,
        thm:peps_dyadic_temporary_holes, lem:peps_dyadic_active_labels,
        thm:peps_protocol_level_events, thm:peps_protocol_party_map_switch,
        thm:peps_dyadic_protocol_parties, thm:peps_dyadic_protocol_final_step,
        thm:peps_guide_frame_norm, thm:peps_guide_frame_supplied_gates,
        lem:peps_dyadic_sampled_guides, thm:peps_protocol_level_frame_facts,
        def:peps_protocol_level_frame, def:peps_dyadic_block_guide,
        thm:peps_dyadic_repainting_valid_labels, def:peps_guide_frame,
        def:peps_dyadic_protocol_schedule, def:peps_dyadic_used_parties,
        lem:peps_adaptive_patch_supply, lem:peps_level_true_vertices,
        thm:peps_protocol_bounded_changes, lem:peps_protocol_junction_resize_parties,
        def:peps_frame}
    Rank monotonicity and coverage follow from Theorem~\ref{thm:peps_dyadic_protocol_schedule_counts}, and when the rank increases the completed block guide of level $j$ equals the starting guide of level $j-1$ (Lemma~\ref{lem:peps_dyadic_active_labels}). Frame congruence (Theorem~\ref{thm:peps_protocol_level_frame_facts}) identifies the transitions across rank boundaries. The root guide has no true vertices (Theorem~\ref{thm:peps_protocol_level_events}) and cell centers lie in $(0,2^k)^2$ (Theorem~\ref{thm:peps_protocol_party_map_switch}), giving the initial raw frame; at the unit level, $2h_1\le 2\epsilon_0<1/2$ gives empty outer physical squares by Theorems~\ref{thm:peps_guide_frame_norm} and~\ref{thm:peps_dyadic_protocol_final_step}.
\end{proof}

\begin{theorem}[The dyadic protocol from its large-scale repaintings]
    \label{thm:peps_protocol_from_large_repaintings}
    \lean{TNLean.PEPS.Approximation.auxIndices,
        TNLean.PEPS.Approximation.exists_large_of_mem_auxIndices,
        TNLean.PEPS.Approximation.mem_eventSchedule_of_mem_auxIndices,
        TNLean.PEPS.Approximation.auxState,
        TNLean.PEPS.Approximation.auxState_succ,
        TNLean.PEPS.Approximation.largeSquare,
        TNLean.PEPS.Approximation.protocolState,
        TNLean.PEPS.Approximation.IsProtocolConstants.t_pos,
        TNLean.PEPS.Approximation.IsProtocolConstants.hole_le,
        TNLean.PEPS.Approximation.exists_isDistributionProtocol_of_largeRepaintings_of_le,
        TNLean.PEPS.Approximation.exists_isDistributionProtocol_of_largeRepaintings}
    \leanok
    \uses{def:peps_dyadic_protocol, def:peps_adaptive_patch_constraint,
        thm:peps_protocol_constants, thm:peps_dyadic_direct_repainting_constants,
        thm:peps_protocol_main_frame_schedule, def:peps_dyadic_protocol_schedule,
        thm:peps_protocol_event_segments, def:peps_circuit_source_bound,
        def:peps_dyadic_protocol_parties, def:peps_frame, thm:peps_protocol_level_events}
    Fix $0<\epsilon_0\le1/16$ and let $\nu'=\min(\epsilon_0,1)/8$. For every
    $0<\nu\le\nu'$, every $k\in\N$, constants of the protocol with an
    admissible direct-repainting ratio, $q\ge1$, $L\ge2$ with $L\le2^k<2L$,
    $\lVert\Omega\rVert=1$ and patches supplied by the adaptive patch constraint
    with constants $B_0,C,c$ at accuracy $L^{-60}$, the following holds. Put
    $t=D\log L$ and let $b$ be the patch bound of
    Theorem~\ref{thm:peps_protocol_constants}. The protocol constants and $L\ge2$
    give $t>0$ and $\epsilon_0t\le B_0\log L$. Let $C',c',k'\in\N$ with
    $C\le C'$, $1\le C'$ and $c\le c'$, and set
    \begin{align}
        K & =(C'+1)^{3b}\bigl(16b^2(C'+1)^{6b}+1\bigr)^{2b}
        L^{3bc'+4b(3bc'+30)}.
        \label{eq:peps_protocol_rewrite_bound}
    \end{align}
    Let the memory before each event be the main frame followed by one frame for
    each auxiliary sheet of the schedule, frozen if its repainting is done and
    an idle raw copy held by its future owner otherwise. Each auxiliary index
    belongs to a scheduled large-scale repainting, and an event other than a
    large-scale repainting freezes no auxiliary sheet. Suppose that each
    repainting of a square of side larger than $K_0t$ passes the memory from its
    state before to its state after through $189$ bounded changes, each with at
    most $2b$ pair effects and $2b$ pair sources in each monomial, at most $K$
    monomials and coefficient sum, error at most $(2b+1)L^{-20}$, and
    participants used by the repainting. Suppose also that the frozen frames
    satisfy the frame constraint at accuracy $L^{-60}$ and together carry at
    most $C_EL^3$ holes. If $5042\le C_N$, $K\le C_KL^{k'}$ and $2b+1\le C_E$,
    there is a protocol with the assembly bounds of
    Definition~\ref{def:peps_dyadic_protocol} for the count constant $C_N$, the
    expansion constant $C_K$, the error constant $C_E$, and the constants
    $64\cdot189$, $k'$ and $2b$, including the bound $2b$ on the pair sources
    of each monomial. This is not
    Proposition~\ref{prop:peps_dyadic_protocol}: the large-scale repaintings are
    assumed here, together with the frozen-frame bounds
    (\texttt{06-geometry.tex}, lines 82--241, 507--541 and 572--681).
\end{theorem}

\begin{proof}\leanok
    \uses{thm:peps_protocol_main_frame_schedule, thm:peps_protocol_level_events,
        thm:peps_protocol_event_segments, thm:peps_dyadic_protocol_schedule_counts,
        def:peps_dyadic_protocol_schedule, def:peps_frame, thm:peps_protocol_constants}
    By Theorem~\ref{thm:peps_protocol_main_frame_schedule}, the main frame follows the schedule transitions from the initial raw copy to the final unit-square frame. Resizes and direct repaintings ($2^j\le K_0t$) act on the main frame by a single rewrite of Theorem~\ref{thm:peps_protocol_level_events} and freeze no auxiliary sheet, while large repaintings ($2^j>K_0t$) use the assumed $189$-step segments (matching the operation count of Theorem~\ref{thm:peps_dyadic_protocol_schedule_counts}). Theorem~\ref{thm:peps_protocol_event_segments} joins all segments into the protocol.
\end{proof}

\begin{theorem}[The side of a non-root dyadic square]
    \label{thm:peps_repainted_square_side_bound}
    \lean{TNLean.PEPS.Approximation.RepaintedSquare.side_le}
    \leanok
    \uses{def:peps_dyadic_protocol_parties}
    If $2^k<2L$, every non-root square of level $j<k$ has side $2^j\le L$.
\end{theorem}

\begin{proof}\leanok
    Since $j+1\le k$, we have $2\cdot2^j=2^{j+1}\le2^k<2L$, so $2^j\le L$.
\end{proof}

\begin{definition}[The frame of a treated sheet]
    \label{def:peps_protocol_treated_frame}
    \lean{TNLean.PEPS.Approximation.treatedGuide,
        TNLean.PEPS.Approximation.treatedVertices,
        TNLean.PEPS.Approximation.treatedFrame,
        TNLean.PEPS.Approximation.treatedFrame_eq}
    \leanok
    \uses{def:peps_adaptive_patch_constraint,
        def:peps_dyadic_block_guide,
        def:peps_dyadic_protocol_parties,
        def:peps_frame,
        def:peps_frame_square_patch,
        def:peps_guide_frame,
        def:peps_protocol_level_frame,
        def:polynomial_peps_model,
        lem:peps_adaptive_patch_supply,
        lem:peps_dyadic_floor,
        lem:peps_treated_true_vertices,
        thm:peps_guide_frame_supplied_gates,
        thm:peps_protocol_bounded_changes}
    Given level $j\in\N$, square index $S\in\Z^2$, $t\in\R$, a finite mark set $M$ of grid corners, homogenizing label $L_0$, and a local guide $g$, the treated guide places $g$ on the block $nS+[0,n]^2$ ($n=2^j$), homogenizes it to $L_0$ on $\bigcup_{Q\in M}B_\infty(nQ,t)$, and samples along $(1,1)$ as in Definition~\ref{def:peps_protocol_level_frame}. Its treated vertex set is the set of true vertices among the grid corners in $[0,N+1]^2$ ($N=2^{k-j}$) and the finite rim patterns of radius $t$ about $M$.

    Given an adaptive patch supply $P$ for $\Omega$ with parameters $(B_0,C,c,60,L)$, a label rank, a repaintable square $R$ at level $j<k$, and $0<\epsilon_0t\le B_0\log L$, the treated frame $F_M^{L_0}(g)$ is the encoded guide frame (Definition~\ref{def:peps_guide_frame}) with supplied patches of radius $\epsilon_0t$ at the treated vertex set when their outer physical samples are pairwise disjoint, and the raw frame owned by $R$ otherwise.
\end{definition}

\begin{proof}\leanok
    Pairwise disjointness of outer physical samples selects the encoded guide-frame branch by construction.
\end{proof}

\begin{definition}[A bounded change of one frame]
    \label{def:peps_frame_change}
    \lean{TNLean.PEPS.Approximation.HasFrameChange}
    \leanok
    \uses{def:peps_bounded_change,
        def:peps_circuit_source_bound,
        def:peps_frame,
        def:peps_frame_registers,
        def:peps_original_circuit,
        def:peps_sheet_stack}
    For $q\ge1$, a bounded frame change between frames $F,F'$ with participant set $S$ and bounds $(r,K,r_s,\eta)$ is an original circuit between their register layouts that is a bounded change (Definition~\ref{def:peps_bounded_change}) from $\Omega_F$ to $\Omega_{F'}$ with parameters $(S,r,K,\eta)$ and has $\le r_s$ pair sources per expansion monomial (Definition~\ref{def:peps_circuit_source_bound}).
\end{definition}

\begin{theorem}[The treated hole radius]
    \label{thm:peps_treated_radius_bounds}
    \lean{TNLean.PEPS.Approximation.IsProtocolConstants.treatedRadius_bounds}
    \leanok
    \uses{thm:peps_protocol_constants}
    For any $L\ge2$ and protocol constants $(K_0,\epsilon_0,\nu,a,D,B_0)$ (Theorem~\ref{thm:peps_protocol_constants}), $0<\epsilon_0D\log L\le B_0\log L$ (\texttt{06-geometry.tex}, lines 564--567).
\end{theorem}

\begin{proof}\leanok
    \uses{thm:peps_protocol_from_large_repaintings}
    Immediate from $\epsilon_0>0$, $D\log L>0$, and $\epsilon_0D\le K_0D\le B_0$ (Theorem~\ref{thm:peps_protocol_from_large_repaintings}).
\end{proof}

\begin{theorem}[The frame of the sampled treated guide]
    \label{thm:peps_treated_frame_spec}
    \lean{TNLean.PEPS.Approximation.treatedFrame_spec}
    \leanok
    \uses{def:peps_adaptive_patch_constraint,
        def:peps_dyadic_block_guide,
        def:peps_dyadic_protocol_parties,
        def:peps_dyadic_used_parties,
        def:peps_guide_frame,
        def:peps_protocol_treated_frame,
        def:polynomial_peps_model,
        lem:peps_adaptive_patch_supply,
        lem:peps_dyadic_birth_enlargement,
        lem:peps_dyadic_floor,
        lem:peps_dyadic_sampled_guides,
        lem:peps_dyadic_unmodified_vertices,
        thm:peps_dyadic_repainting_valid_labels,
        thm:peps_guide_frame_supplied_gates,
        thm:peps_protocol_bounded_changes,
        thm:peps_protocol_constants,
        thm:peps_treated_radius_bounds}
    Under protocol constants with $L\ge2$, $t=D\log L$, a repaintable square $R$ at level $j<k$ of side $2^j>K_0t$, an adaptive patch supply $P$ with parameters $(B_0,C,c,60,L)$, and any unmodified guide $g$ of the repainting baseline of $R$, finite mark set $M$, and label $L_0$, the treated guide of Definition~\ref{def:peps_protocol_treated_frame} is a $(1,1)$-reading of the homogenized guide on $R$, its true vertices equal the treated vertex set, their supplied outer samples are pairwise disjoint, and $F_M^{L_0}(g)$ equals the encoded guide frame (\texttt{06-geometry.tex}, lines 69--80, 102--125 and 328--355).
\end{theorem}

\begin{proof}\leanok
    \uses{def:peps_dyadic_block_guide,
        def:peps_protocol_treated_frame,
        lem:peps_dyadic_chamber_sampling,
        lem:peps_dyadic_readings_exist,
        lem:peps_treated_true_vertices,
        thm:peps_dyadic_repainting_valid_labels,
        thm:peps_dyadic_temporary_holes,
        thm:peps_guide_frame_norm,
        thm:peps_guide_frame_supplied_gates,
        thm:peps_protocol_bounded_changes,
        thm:peps_protocol_level_frame_facts,
        thm:peps_treated_radius_bounds}
    Ray-constancy (Lemma~\ref{lem:peps_dyadic_readings_exist}) gives the $(1,1)$-reading, Lemma~\ref{lem:peps_treated_true_vertices} and Theorem~\ref{thm:peps_dyadic_temporary_holes} identify the true vertices and their separation $>4\epsilon_0t$, and disjointness selects the encoded guide-frame branch of Definition~\ref{def:peps_protocol_treated_frame}.
\end{proof}

\begin{theorem}[A lower bound for the protocol patch count]
    \label{thm:peps_protocol_patch_bound_lower}
    \lean{TNLean.PEPS.Approximation.fiftyThree_le_protocolPatchBound}
    \leanok
    \uses{thm:peps_protocol_constants}
    For $\epsilon_0\ge0$ and $0<\nu<1/3$, the protocol patch bound $b$ of Theorem~\ref{thm:peps_protocol_constants} satisfies $b\ge53$.
\end{theorem}

\begin{proof}\leanok
    Since $0<\nu<1/3$ and $\epsilon_0\ge0$, we have $(4+8\epsilon_0+2\nu)/\nu+1\ge 8$, so $b\ge 8^2=64\ge 53$.
\end{proof}

\begin{theorem}[A point treatment as a change of a treated sheet]
    \label{thm:peps_protocol_point_treatment_change}
    \lean{TNLean.PEPS.Approximation.pointTreatment_hasSourceBoundedChange_of_le,
        TNLean.PEPS.Approximation.exists_pointTreatment_hasSourceBoundedChange}
    \leanok
    \uses{def:peps_adaptive_patch_constraint,
        def:peps_dyadic_block_guide,
        def:peps_dyadic_protocol_parties,
        def:peps_dyadic_used_parties,
        def:peps_frame_change,
        def:peps_protocol_treated_frame,
        def:polynomial_peps_model,
        lem:peps_adaptive_patch_supply,
        lem:peps_dyadic_direct_repainting,
        lem:peps_dyadic_unmodified_vertices,
        thm:peps_dyadic_block_ratio,
        thm:peps_dyadic_repainting_valid_labels,
        thm:peps_protocol_bounded_changes,
        thm:peps_protocol_constants,
        thm:peps_protocol_level_events,
        thm:peps_treated_radius_bounds}
    Fix $0<\epsilon_0\le1/16$ and let $\nu'=\min(\epsilon_0,1)/20000$. The following holds for
    every $0<\nu\le\nu'$. Choose protocol constants, natural
    $L\ge2$ and $q\ge1$, a unit vector $\Omega$, an adaptive supplied-patch family $P$ for
    $\Omega$ with parameters $B_0,C,c$, test exponent $60$ and size $L$, and any label rank.
    Let $R$ be a repaintable square of a hierarchy of any natural depth $k$, at level $j<k$,
    with integer index $s$, side $n=2^j>K_0t$ and $t=D\log L$. Choose any set of already
    repainted blocks and an unmodified guide $g$ of its repainting baseline. Let $L_0$ be the
    new label of $R$ or the current label of one of the nine blocks with indices $s+u$,
    $u\in\{-1,0,1\}^2$. Let $M,M'$ be finite sets of grid corners with identical membership
    away from one corner $Q_0$ of $R$; equality of the two sets is allowed.

    For natural $C',c'$ satisfying $C\le C'$, $1\le C'$ and $c\le c'$, put $b$ equal to the
    protocol patch bound and use the bound $K_{\mathrm{rw}}$
    of~(\ref{eq:peps_protocol_level_rewrite_bound}) for these $C',c',b,L$. Passing from the
    treated frame with marks $M$ to that with marks $M'$ is a bounded change with participants
    used by the repainting of $R$, pair-effect and pair-source bounds $2b$, monomial bound
    $K_{\mathrm{rw}}$, and error $(2b+1)L^{-20}$ (\texttt{06-geometry.tex}, lines 328--339,
    481--505 and 572--596).
\end{theorem}

\begin{proof}\leanok
    \uses{def:peps_dyadic_chamber_sampling,
        lem:peps_treated_true_vertices,
        thm:peps_guide_frame_supplied_gates,
        thm:peps_protocol_bounded_changes,
        thm:peps_protocol_patch_bound_lower,
        thm:peps_treated_frame_spec,
        thm:peps_treated_radius_bounds}
    The $(1,1)$-readings of the two homogenized guides agree outside $\overline B_\infty(Q_0,t)$ and their frames equal the encoded guide frames by Theorem~\ref{thm:peps_treated_frame_spec}. Rewriting the supplied radius-$\epsilon_0t$ patches via auxiliary patches of radius $u=\nu t/2\le B_0\log L$ (Theorem~\ref{thm:peps_guide_frame_supplied_gates}) uses a square net of size $\le((4+8\epsilon_0+2\nu)/\nu+1)^2\le b$ with $b\ge53$ (Theorem~\ref{thm:peps_protocol_patch_bound_lower}), yielding the bounds $(2b,K_{\mathrm{rw}},2b,(2b+1)L^{-20})$ and participants used by $R$ (Theorem~\ref{thm:peps_protocol_bounded_changes}).
\end{proof}

\begin{theorem}[Readings separated from a change]
    \label{thm:peps_readings_far_change}
    \lean{TNLean.PEPS.Approximation.readings_isTrueVertex_iff_of_far}
    \leanok
    \uses{def:peps_dyadic_incident_labels,
        lem:peps_dyadic_birth_enlargement,
        lem:peps_dyadic_sampled_guides}
    If plane labellings $G_b,G_a$ have the same true vertices, each at distance $\ge\delta>0$ from $\overline{\{G_b\ne G_a\}}$, then their readings $f_b,f_a$ along any common displacement have the same true vertices and the same incident labels at them.
\end{theorem}

\begin{proof}\leanok
    \uses{def:peps_dyadic_incident_labels,
        lem:peps_dyadic_incident_change,
        lem:peps_dyadic_sampled_guides,
        thm:peps_dyadic_temporary_holes}
    Since $\{f_b\ne f_a\}\subseteq\overline{\{G_b\ne G_a\}}$ (Lemma~\ref{lem:peps_dyadic_sampled_guides}), $f_b$ and $f_a$ agree on a neighborhood of every guide vertex, so their incident labels and true-vertex sets coincide (Theorem~\ref{thm:peps_dyadic_temporary_holes}).
\end{proof}

\begin{theorem}[Changing the owners of a guide frame]
    \label{thm:peps_guide_frame_change_owner}
    \lean{TNLean.PEPS.Approximation.guideFrame_eq_changeOwner,
        TNLean.PEPS.Approximation.guideFrame_eq_changeOwner_of_ne}
    \leanok
    \uses{def:peps_dyadic_incident_labels,
        def:peps_frame_birth_partition,
        def:peps_frame_square_patch,
        def:peps_guide_frame}
    For two guides $g,g'$ sharing a finite center set $V$, patch family with disjoint outer samples, owner map, placeholder, and label rank, if $g'$ assigns owner $Q$ on a site set $T$, agrees with the owners of $g$ off $T$, and has the same tag holders on $V$, then the guide frame of $g'$ is $F[T\mapsto Q]$ (where $F$ is the guide frame of $g$). In particular, if $T=\{s\mid g(s)\ne g'(s)\}$, $g'\equiv B$ on $T$, and $g,g'$ have equal incident label sets at every $z\in V$, then the guide frame of $g'$ is $F[T\mapsto\text{owner}(B)]$.
\end{theorem}

\begin{proof}\leanok
    \uses{def:peps_guide_frame,
        lem:peps_guide_frame_tags,
        thm:peps_protocol_main_frame_schedule}
    Site owners match $F[T\mapsto Q]$ everywhere and the ordered tagged holes on $V$ are identical (Definition~\ref{def:peps_guide_frame}); in the corollary, equal incident label sets on $V$ yield equal canonical tag holders (Lemma~\ref{lem:peps_guide_frame_tags}).
\end{proof}

\begin{theorem}[Values where two readings differ]
    \label{thm:peps_reading_changed_values}
    \lean{TNLean.PEPS.Approximation.IsDisplacedReading.values_of_ne}
    \leanok
    \uses{def:peps_dyadic_block_guide,
        lem:peps_dyadic_sampled_guides}
    If labellings $H_b,H_a$ satisfy $H_b=A$ and $H_a\in\{A,B\}$ on the open $\delta$-neighborhood ($\delta>0$) of $\{H_b\ne H_a\}$, then for any $(1,1)$-readings $f_b,f_a$ of their translates by $d$, every point with $f_b(p)\ne f_a(p)$ satisfies $f_b(p)=A$ and $f_a(p)=B$.
\end{theorem}

\begin{proof}\leanok
    \uses{lem:peps_dyadic_sampled_guides,
        lem:peps_dyadic_translation}
    Any $p$ with $f_b(p)\ne f_a(p)$ has $p-d\in\overline{\{H_b\ne H_a\}}$ (Lemmas~\ref{lem:peps_dyadic_sampled_guides} and~\ref{lem:peps_dyadic_translation}), so for small $\eta>0$, $p-d+\eta(1,1)$ lies in the $\delta$-neighborhood where $f_b(p)=A$ and $f_a(p)\in\{A,B\}\setminus\{A\}=\{B\}$.
\end{proof}

\begin{theorem}[Support of a treated band change]
    \label{thm:peps_treated_band_support}
    \lean{TNLean.PEPS.Approximation.BandOperation.treated_ne_subset}
    \leanok
    \uses{def:peps_dyadic_edge_band,
        def:peps_dyadic_normal_words,
        def:peps_dyadic_schedule,
        lem:peps_dyadic_floor}
    For any edge band operation, real $n,t$, edge, and background guide, the difference set of its endpoint-treated before and after guides is contained in the open edge band $\{-8<x<2\}$.
\end{theorem}

\begin{proof}\leanok
    The two endpoint-treated guides agree on the endpoint squares and both equal the background guide outside the open band $\{-8<x<2\}$.
\end{proof}

\begin{theorem}[Values near a treated band change]
    \label{thm:peps_treated_band_values}
    \lean{TNLean.PEPS.Approximation.BandOperation.treated_mem}
    \leanok
    \uses{def:peps_dyadic_edge_band,
        def:peps_dyadic_normal_words,
        def:peps_dyadic_schedule,
        lem:peps_dyadic_floor,
        thm:peps_dyadic_no_changed_vertex}
    For an edge band operation inserting a single label $Q$ with surrounding label $P_\circ$ and $t\le n$, on the open $a_0t$-neighborhood of the difference of its endpoint-treated guides both treated guides take values in $\{P_\circ,Q\}$ and the treated surrounding guide equals $P_\circ$.
\end{theorem}

\begin{proof}\leanok
    \uses{thm:peps_dyadic_no_changed_vertex}
    By Theorem~\ref{thm:peps_dyadic_no_changed_vertex}, the treated surrounding guide is $P_\circ$ on this neighborhood, and the single-label insertion property restricts both endpoint-treated guides to $\{P_\circ,Q\}$.
\end{proof}

\begin{theorem}[Values of the treated band readings]
    \label{thm:peps_treated_band_readings}
    \lean{TNLean.PEPS.Approximation.BandOperation.readings_of_ne}
    \leanok
    \uses{def:peps_dyadic_block_guide,
        def:peps_dyadic_edge_band,
        def:peps_dyadic_normal_words,
        def:peps_dyadic_protocol_parties,
        def:peps_dyadic_schedule,
        lem:peps_dyadic_floor,
        lem:peps_dyadic_sampled_guides,
        thm:peps_dyadic_no_changed_vertex}
    For an elementary edge operation inserting $Q$ with surrounding label $P_\circ$ and $0<t\le n$, let $f_b,f_a,f_\circ$ be $(1,1)$-readings of translates of the endpoint-treated before, after, and surrounding guides. At every point where $f_b\ne f_a$, one has $f_\circ=P_\circ$ and $f_b,f_a\in\{P_\circ,Q\}$.
\end{theorem}

\begin{proof}\leanok
    \uses{lem:peps_dyadic_sampled_guides,
        lem:peps_dyadic_translation,
        thm:peps_treated_band_values}
    A point where $f_b\ne f_a$ translates into the closure of the local guide difference, whose open $a_0t$-neighborhood contains small positive $(1,1)$-displacements; Theorem~\ref{thm:peps_treated_band_values} concludes.
\end{proof}

\begin{theorem}[Geometry of the treated band readings]
    \label{thm:peps_treated_band_frame_facts}
    \lean{TNLean.PEPS.Approximation.BandOperation.treated_frame_facts}
    \leanok
    \uses{def:peps_dyadic_block_guide,
        def:peps_dyadic_edge_band,
        def:peps_dyadic_incident_labels,
        def:peps_dyadic_normal_words,
        def:peps_dyadic_protocol_parties,
        def:peps_dyadic_schedule,
        lem:peps_dyadic_birth_enlargement,
        lem:peps_dyadic_floor,
        lem:peps_dyadic_sampled_guides,
        thm:peps_dyadic_no_changed_vertex}
    Under the hypotheses of Theorem~\ref{thm:peps_treated_band_readings} with $2\epsilon_0<a_0$, the $(1,1)$-readings $f_b,f_a$ have the same true vertices and incident labels, every $p\in\{f_b\ne f_a\}$ satisfies $d_\infty(p,z)>2\epsilon_0t$ for every true vertex $z$, and $d_\infty(p,p')\le n$ for all $p,p'\in\{f_b\ne f_a\}$.
\end{theorem}

\begin{proof}\leanok
    \uses{lem:peps_dyadic_band_disjoint,
        lem:peps_dyadic_sampled_guides,
        lem:peps_dyadic_translation,
        thm:peps_bulk_steps_treated_frames,
        thm:peps_dyadic_no_changed_vertex,
        thm:peps_dyadic_temporary_holes,
        thm:peps_readings_far_change,
        thm:peps_treated_band_support}
    By Theorem~\ref{thm:peps_dyadic_no_changed_vertex}, the treated guides share the same true vertices at distance $\ge a_0t>2\epsilon_0t$ from the closure of their difference, so Theorem~\ref{thm:peps_readings_far_change} gives equal reading vertices, incident labels, and clearance $>2\epsilon_0t$. Theorem~\ref{thm:peps_treated_band_support} places $\{f_b\ne f_a\}$ inside the closed edge band of diameter $\le n$ (Theorem~\ref{thm:peps_bulk_steps_treated_frames}).
\end{proof}

\begin{theorem}[A birth along an edge as a change of a treated sheet]
    \label{thm:peps_protocol_band_step_changes}
    \lean{TNLean.PEPS.Approximation.bandBirth_hasFrameChange}
    \leanok
    \uses{def:peps_adaptive_patch_constraint,
        def:peps_angular_scale,
        def:peps_dyadic_block_guide,
        def:peps_dyadic_edge_band,
        def:peps_dyadic_normal_words,
        def:peps_dyadic_protocol_parties,
        def:peps_dyadic_schedule,
        def:peps_dyadic_used_parties,
        def:peps_frame_change,
        def:peps_protocol_treated_frame,
        def:polynomial_peps_model,
        lem:peps_adaptive_patch_supply,
        lem:peps_dyadic_unmodified_vertices,
        thm:peps_bulk_steps_treated_frames,
        thm:peps_dyadic_no_changed_vertex,
        thm:peps_dyadic_repainting_valid_labels,
        thm:peps_protocol_bounded_changes,
        thm:peps_protocol_constants,
        thm:peps_treated_radius_bounds}
    Let $q\ge1$, fix a family of states, and let $\Omega$ be a unit vector of the family.
    Assume that $D$ is an angular scale as in Definition~\ref{def:peps_angular_scale} for
    the family with parameters $q,a,A_0,m_0$, test
    exponent $60$ and threshold $L_0'$, where $m_0\ge2$ and $A_0\ge1$. Choose protocol
    constants and natural $L\ge\max(2,L_0')$ and $k$ with $2^k<2L$. Fix an adaptive
    supplied-patch family $P$ for $\Omega$ with parameters $B_0,C,c$, test exponent $60$ and
    size $L$, and any label rank.

    Let $R$ be a repaintable square at level $j<k$, with $n=2^j>K_0t$ and $t=D\log L$, and
    choose any set of already repainted blocks. An elementary edge birth involving a single
    label $Q$ changes one unmodified guide of this repainting baseline into another. Its
    surrounding guide is the guide before. Treat both guides with the surrounding label about
    exactly the endpoints of the edge. If the parties of the surrounding label and of $Q$ lie
    in a participant set $S$, passing from the treated frame before to that after is a bounded
    change with participants in $S$, at most one pair source and error $L^{-30}$. This holds
    for all natural monomial and pair-source bounds $K,r_s\ge1$; the natural pair-effect bound
    $r$ is arbitrary (\texttt{06-geometry.tex}, lines 183--221, 328--387 and 440--445).
\end{theorem}

\begin{proof}\leanok
    \uses{def:peps_frame_change,
        lem:peps_guide_frame_outer_holes_disjoint,
        lem:peps_adaptive_patch_supply,
        lem:peps_circuit_source_bound_monomials,
        thm:peps_bounded_change_closure,
        thm:peps_bulk_steps_treated_frames,
        thm:peps_guide_frame_change_owner,
        thm:peps_guide_frame_supplied_gates,
        thm:peps_repainted_square_side_bound,
        thm:peps_treated_band_frame_facts,
        thm:peps_treated_band_readings,
        thm:peps_treated_frame_spec,
        thm:peps_treated_radius_bounds}
    By Theorems~\ref{thm:peps_treated_frame_spec}, \ref{thm:peps_treated_band_readings}, and~\ref{thm:peps_treated_band_frame_facts}, the before and after readings share the same true vertices and incident labels, and the changed site set $T$ avoids all outer holes ($d_\infty(p,z)>2\epsilon_0t$), has diameter $\le n\le L$ (Theorem~\ref{thm:peps_repainted_square_side_bound}), and switches labels from $P_\circ$ to $Q$. Thus $F_a=F_b[T\mapsto Q_\circ]$ by Theorem~\ref{thm:peps_guide_frame_change_owner}, and the bulk birth of Theorem~\ref{thm:peps_bulk_steps_treated_frames} (converted via Theorem~\ref{thm:peps_bounded_change_closure}) gives the bounded frame change with $\le1$ pair source and error $L^{-30}$.
\end{proof}

\begin{theorem}[A death along an edge as a change of a treated sheet]
    \label{thm:peps_protocol_band_death_change}
    \lean{TNLean.PEPS.Approximation.bandDeath_hasFrameChange}
    \leanok
    \uses{def:peps_adaptive_patch_constraint,
        def:peps_angular_scale,
        def:peps_dyadic_block_guide,
        def:peps_dyadic_edge_band,
        def:peps_dyadic_normal_words,
        def:peps_dyadic_protocol_parties,
        def:peps_dyadic_schedule,
        def:peps_dyadic_used_parties,
        def:peps_frame_change,
        def:peps_protocol_treated_frame,
        def:polynomial_peps_model,
        lem:peps_adaptive_patch_supply,
        lem:peps_dyadic_unmodified_vertices,
        thm:peps_bulk_steps_treated_frames,
        thm:peps_dyadic_no_changed_vertex,
        thm:peps_dyadic_repainting_valid_labels,
        thm:peps_protocol_bounded_changes,
        thm:peps_protocol_constants,
        thm:peps_treated_radius_bounds}
    Let $q\ge1$, fix a family of states, and let $\Omega$ be a unit vector of the family.
    Assume that $D$ is an angular scale as in Definition~\ref{def:peps_angular_scale} for
    the family with parameters $q,a,A_0,m_0$, test
    exponent $60$ and threshold $L_0'$, where $m_0\ge2$ and $A_0\ge1$. Choose protocol
    constants and natural $L\ge\max(2,L_0')$ and $k$ with $2^k<2L$. Fix an adaptive
    supplied-patch family $P$ for $\Omega$ with parameters $B_0,C,c$, test exponent $60$ and
    size $L$, and any label rank.

    Let $R$ be a repaintable square at level $j<k$, with $n=2^j>K_0t$ and $t=D\log L$, and
    choose any set of already repainted blocks. An elementary edge death involving a single
    label $Q$ changes one unmodified guide of this repainting baseline into another. Its
    surrounding guide is the guide after. Treat both guides with the surrounding label about
    exactly the endpoints of the edge. If the parties of the surrounding label and of $Q$ lie
    in a participant set $S$, passing from the treated frame before to that after is a bounded
    change with participants in $S$, at most one pair source and error $L^{-30}$. This holds
    for all natural monomial and pair-source bounds $K,r_s\ge1$; the natural pair-effect bound
    $r$ must satisfy $r\ge1$ (\texttt{06-geometry.tex}, lines 183--221, 328--387 and
    440--445).
\end{theorem}

\begin{proof}\leanok
    \uses{def:peps_frame_change,
        lem:peps_guide_frame_outer_holes_disjoint,
        lem:peps_adaptive_patch_supply,
        lem:peps_circuit_source_bound_monomials,
        thm:peps_bounded_change_closure,
        thm:peps_bulk_steps_treated_frames,
        thm:peps_guide_frame_change_owner,
        thm:peps_guide_frame_supplied_gates,
        thm:peps_repainted_square_side_bound,
        thm:peps_treated_band_frame_facts,
        thm:peps_treated_band_readings,
        thm:peps_treated_frame_spec,
        thm:peps_treated_radius_bounds}
    Just as in Theorem~\ref{thm:peps_protocol_band_step_changes}, Theorem~\ref{thm:peps_guide_frame_change_owner} gives $F_b=F_a[T\mapsto Q_\circ]$, and the bulk death of Theorem~\ref{thm:peps_bulk_steps_treated_frames} (which uses one pair effect, requiring $r\ge1$) yields the bounded frame change with error $L^{-30}$.
\end{proof}

\begin{theorem}[The support of a treated central birth]
    \label{thm:peps_central_treated_support}
    \lean{TNLean.PEPS.Approximation.RepaintingBaseline.central_treated_ne_subset}
    \leanok
    \uses{def:peps_dyadic_schedule,
        lem:peps_dyadic_floor}
    For any repainting baseline of side $n$ and real $t$, homogenizing its starting and central guides to its old label on the open radius-$t$ sup squares about the four corners produces guides that differ only inside $[0,n]^2$.
\end{theorem}

\begin{proof}\leanok
    The two guides agree on the four corner squares by homogenization, and outside the central region (which lies in $[0,n]^2$) the central guide equals the starting guide.
\end{proof}

\begin{theorem}[The central birth as a change of the main sheet]
    \label{thm:peps_protocol_central_birth_change}
    \lean{TNLean.PEPS.Approximation.centralBirth_hasFrameChange_of,
        TNLean.PEPS.Approximation.centralBirth_hasFrameChange}
    \leanok
    \uses{def:peps_adaptive_patch_constraint,
        def:peps_angular_scale,
        def:peps_dyadic_block_guide,
        def:peps_dyadic_protocol_parties,
        def:peps_dyadic_schedule,
        def:peps_dyadic_used_parties,
        def:peps_frame_change,
        def:peps_frame_square_patch,
        def:peps_guide_frame,
        def:peps_protocol_treated_frame,
        def:polynomial_peps_model,
        lem:peps_adaptive_patch_supply,
        lem:peps_dyadic_birth_enlargement,
        lem:peps_dyadic_floor,
        lem:peps_dyadic_sampled_guides,
        thm:peps_bulk_steps_treated_frames,
        thm:peps_dyadic_repainting_valid_labels,
        thm:peps_guide_frame_supplied_gates,
        thm:peps_protocol_bounded_changes,
        thm:peps_protocol_constants,
        thm:peps_treated_radius_bounds}
    For the sampled-guide change, let $k,L,q$ be natural numbers with $q\ge1$,
    fix a family of states, and let $\Omega$ be a unit vector of the family on
    the $L\times L$ system. Assume the exponent-$60$ angular-scale condition
    for the family with parameters $q,a,A_0,m_0,D,L'_0$, with natural $m_0$, where
    $L'_0\le L$, $m_0\ge4$, $A_0\ge1$ and $a\le a_0/2$.
    Here $a_0$ is the angular constant. Take any repainting baseline on the
    protocol guide labels, any translation $d$, and real $t,\epsilon_0$ with
    $0<t=D\log L\le n\le A_0L$ and $2\epsilon_0\le a_0/2$.
    Let $f_b,f_a$ be the $(1,1)$ readings of the translates by $d$ of the
    starting and central guides after homogenizing them to the old label on
    the four local corner squares of radius $t$. Let $V_b,V_a$ be finite sets
    consisting exactly of their true vertices. Take an adaptive patch supply at the cell centers for
    $\Omega$, with parameters $B_0,C,c$, test exponent $60$ and size $L$,
    and suppose $0<\epsilon_0t\le B_0\log L$. Within each of $V_b,V_a$,
    require the supplied patches of radius $\epsilon_0t$ to have
    pairwise-disjoint outer physical samples. Fix any owner map from guide
    labels to the level-$k$ protocol parties and any natural-valued label rank.
    The supplied guide frames of $f_b,f_a$, with the protocol placeholder,
    then admit a bounded change with error $L^{-30}$ and participants in
    any set containing the owners of the old and final labels, for every
    natural pair-effect bound $r$ and all natural $K,r_s\ge1$. The bounded-change
    condition bounds both the number of monomials and the sum of absolute
    coefficients by $K$; at most one pair source is needed.

    For the protocol specialization, assume the same family, unit state and
    exponent-$60$ angular-scale data with $q\ge1$,
    $L'_0\le L$, $m_0\ge4$ and $A_0\ge1$. Assume the protocol constants
    $K_0,\epsilon_0,\nu,a,D,B_0$, $L\ge2$ and $2^k<2L$, and take the
    exponent-$60$ adaptive supply and any label rank as above. Let $R$ be
    any repaintable square of level $j<k$, index $S_0$ and side
    $n=2^j>K_0t$, where $t=D\log L$. For any set $\mathcal D$ of block
    indices, form the block baseline at $S_0$ from the level-$j$ protocol
    labels at $\mathcal D$, with final label the new square party of $R$.
    Passing from its starting treated frame to its central treated frame,
    both homogenized to its old label at the four corners of $R$, has the
    same bounds for all natural $r,K,r_s$ with $K,r_s\ge1$, with participants
    in any set containing the parties of its old and final labels. No condition
    $S_0\notin\mathcal D$ is required
    (\texttt{06-geometry.tex}, lines 155--160, 328--387 and 440--445).
\end{theorem}

\begin{proof}\leanok
    \uses{def:peps_frame_change,
        lem:peps_guide_frame_outer_holes_disjoint,
        lem:peps_adaptive_patch_supply,
        lem:peps_circuit_source_bound_monomials,
        lem:peps_dyadic_sampled_guides,
        lem:peps_dyadic_translation,
        thm:peps_bounded_change_closure,
        thm:peps_bulk_steps_treated_frames,
        thm:peps_central_treated_support,
        thm:peps_dyadic_no_changed_vertex,
        thm:peps_dyadic_temporary_holes,
        thm:peps_guide_frame_change_owner,
        thm:peps_guide_frame_supplied_gates,
        thm:peps_reading_changed_values,
        thm:peps_readings_far_change,
        thm:peps_repainted_square_side_bound,
        thm:peps_treated_frame_spec,
        thm:peps_treated_radius_bounds}
    By Theorems~\ref{thm:peps_dyadic_no_changed_vertex}, \ref{thm:peps_readings_far_change}, and~\ref{thm:peps_reading_changed_values}, the two readings share the same true vertices and incident labels, and the changed site set $T$ switches from $A$ to $B$, avoids all outer holes ($d_\infty(p,z)>2\epsilon_0t$), and has diameter $\le n\le L$ (Theorems~\ref{thm:peps_central_treated_support} and~\ref{thm:peps_repainted_square_side_bound}). Theorem~\ref{thm:peps_guide_frame_change_owner} and the bulk birth of Theorem~\ref{thm:peps_bulk_steps_treated_frames} conclude.
\end{proof}

\begin{theorem}[Reading a guide pieced across an open region]
    \label{thm:peps_reading_piecewise}
    \lean{TNLean.PEPS.Approximation.IsDisplacedReading.eq_piecewise}
    \leanok
    \uses{def:peps_dyadic_block_guide,
        lem:peps_dyadic_sampled_guides}
    Let $Y,U\subseteq\R^2$ be open with $\partial Y\subseteq U$, let guides $H_1,H_2$ both equal $C$ on $U$, and let $f_1,f_2,g$ be $(1,1)$-readings of the translates by $d$ of $H_1,H_2$ and of the piecewise guide $H$ equal to $H_2$ on $Y$ and $H_1$ on $Y^c$. Then $g(p)=f_2(p)$ for $p-d\in Y$ and $g(p)=f_1(p)$ for $p-d\notin Y$.
\end{theorem}

\begin{proof}\leanok
    \uses{def:peps_dyadic_block_guide,
        lem:peps_dyadic_sampled_guides}
    Set $u_\eta=p+\eta(1,1)-d$ for small $\eta>0$. Since $Y$, $U$, and $(\overline Y)^c=(Y\cup\partial Y)^c\supseteq(Y\cup U)^c$ are open, $u_\eta$ stays in $Y$ when $p-d\in Y$ (giving $g(p)=H_2(u_\eta)=f_2(p)$), in $U$ when $p-d\in U\setminus Y$ (giving $g(p)=C=H_1(u_\eta)=f_1(p)$), and in $(\overline Y)^c$ when $p-d\notin Y\cup U$ (giving $g(p)=H_1(u_\eta)=f_1(p)$).
\end{proof}

\begin{theorem}[Reading a guide constant on an open region]
    \label{thm:peps_reading_constant_open}
    \lean{TNLean.PEPS.Approximation.IsDisplacedReading.eq_of_mem_open}
    \leanok
    \uses{def:peps_dyadic_block_guide,
        lem:peps_dyadic_sampled_guides}
    If $U\subseteq\R^2$ is open, $H\equiv C$ on $U$, and $f$ is a $(1,1)$-reading of the translate of $H$ by $d$, then $f(p)=C$ for every $p$ with $p-d\in U$.
\end{theorem}

\begin{proof}\leanok
    \uses{lem:peps_dyadic_sampled_guides}
    For small $\eta>0$, $p+\eta(1,1)-d\in U$, so $f(p)=H(p+\eta(1,1)-d)=C$.
\end{proof}

\begin{theorem}[Hole lists of a pieced guide]
    \label{thm:peps_exchange_holes_perm}
    \lean{TNLean.PEPS.Approximation.guideFrame_exchange_holes_perm}
    \leanok
    \uses{def:peps_frame,
        def:peps_frame_square_patch,
        def:peps_guide_frame,
        lem:peps_dyadic_birth_enlargement}
    Let $Y'\subseteq\R^2$, let $U'\subseteq\R^2$ be open with $\partial Y'\subseteq U'$, and suppose guides $f_1,f_2$ both equal $C$ on $U'$. For the piecewise guide $f$ equal to $f_2$ on $Y'$ and $f_1$ on $(Y')^c$, with exact finite true-vertex sets $V_1,V_2,W$ and a common patch family (with pairwise disjoint outer samples on each vertex set), owner map, placeholder, and rank, the holes of $f_1$ centered in $(Y')^c$ followed by the holes of $f_2$ centered in $Y'$ form a permutation of the holes (including tag owners) of the guide frame of $f$.
\end{theorem}

\begin{proof}\leanok
    \uses{def:peps_guide_frame,
        lem:peps_dyadic_incident_change,
        lem:peps_guide_frame_tags}
    Since $f_1=f_2=C$ on $U'\supseteq\partial Y'$, the true vertices $W$ of $f$ partition into $V_1\setminus Y'$ and $V_2\cap Y'$ with matching incident label sets and canonical tag holders (Lemmas~\ref{lem:peps_dyadic_incident_change} and~\ref{lem:peps_guide_frame_tags}).
\end{proof}

\begin{definition}[A bounded change of two frames]
    \label{def:peps_protocol_pair_change}
    \lean{TNLean.PEPS.Approximation.HasPairChange}
    \leanok
    \uses{def:peps_bounded_change,
        def:peps_circuit_source_bound,
        def:peps_frame,
        def:peps_frame_registers,
        def:peps_original_circuit}
    For $q\ge1$, a bounded pair change from $(F_1,F_2)$ to $(F'_1,F'_2)$ with participant set $S$ and bounds $(r,K,r_s,\eta)$ is an original circuit from the concatenated registers of $F_1,F_2$ to those of $F'_1,F'_2$ that is a bounded change (Definition~\ref{def:peps_bounded_change}) from $\Omega_{F_1}\otimes\Omega_{F_2}$ to $\Omega_{F'_1}\otimes\Omega_{F'_2}$ with parameters $(S,r,K,\eta)$ and has $\le r_s$ pair sources per expansion monomial (Definition~\ref{def:peps_circuit_source_bound}).
\end{definition}

\begin{theorem}[Reordering the targets of a bounded pair change]
    \label{thm:peps_pair_change_reorder}
    \lean{TNLean.PEPS.Approximation.HasPairChange.reorder}
    \leanok
    \uses{def:peps_frame,
        def:peps_protocol_pair_change}
    Any bounded pair change (Definition~\ref{def:peps_protocol_pair_change}) into $(N_1,N_2)$ is also a bounded pair change with the same bounds $(S,r,K,r_s,\eta)$ into any $(G_1,G_2)$ having the same site owners and permuted hole lists.
\end{theorem}

\begin{proof}\leanok
    \uses{def:peps_protocol_pair_change,
        lem:peps_circuit_source_bound_monomials,
        lem:peps_register_reorderings,
        thm:peps_bounded_change_closure,
        thm:peps_guide_frame_rewrite}
    By Lemma~\ref{lem:peps_register_reorderings} and Theorem~\ref{thm:peps_guide_frame_rewrite}, tensor-factor reorderings (with $0$ nonprivate gates) map $\Omega_{N_1}\otimes\Omega_{N_2}\mapsto\Omega_{G_1}\otimes\Omega_{G_2}$, preserving all bounds by Theorem~\ref{thm:peps_bounded_change_closure}.
\end{proof}

\begin{theorem}[The lens exchange as a change of two treated sheets]
    \label{thm:peps_protocol_lens_exchange_change}
    \lean{TNLean.PEPS.Approximation.edgeExchange_hasPairChange_of,
        TNLean.PEPS.Approximation.edgeExchange_hasPairChange}
    \leanok
    \uses{def:peps_adaptive_patch_constraint,
        def:peps_angular_scale,
        def:peps_dyadic_block_guide,
        def:peps_dyadic_edge_band,
        def:peps_dyadic_normal_words,
        def:peps_dyadic_protocol_parties,
        def:peps_dyadic_schedule,
        def:peps_dyadic_used_parties,
        def:peps_frame_square_patch,
        def:peps_guide_frame,
        def:peps_protocol_pair_change,
        def:peps_protocol_treated_frame,
        def:polynomial_peps_model,
        lem:peps_adaptive_patch_supply,
        lem:peps_dyadic_birth_enlargement,
        lem:peps_dyadic_floor,
        lem:peps_dyadic_sampled_guides,
        lem:peps_dyadic_unmodified_vertices,
        thm:peps_bulk_steps_treated_frames,
        thm:peps_dyadic_no_changed_vertex,
        thm:peps_dyadic_repainting_valid_labels,
        thm:peps_guide_frame_supplied_gates,
        thm:peps_protocol_bounded_changes,
        thm:peps_protocol_constants,
        thm:peps_treated_radius_bounds}
    For the sampled-guide exchange, let $k,L,q$ be natural with $q\ge1$,
    fix a family of states, and suppose $\Omega$ is a unit vector of the
    family on the $L\times L$ system. Assume the exponent-$60$ angular-scale
    condition for the family with parameters $q,a,A_0,m_0,D,L'_0$, with natural $m_0$, where
    $L'_0\le L$, $m_0\ge2$, $A_0\ge1$ and $a\le a_0/2$, with $a_0$ the angular
    constant. Take any repainting baseline on protocol guide labels, any
    edge list $E$, edge $e$ and translation $d$. Require
    $0<t=D\log L\le n\le A_0L$ and
    $2\epsilon_0(a_0+1)\le a_0/2$.
    Let $C_e$ be the neighboring label along $e$. Take the four $(1,1)$
    readings of the translates by $d$ of the guides after treating them
    with $C_e$ on the radius-$t$ squares about the local endpoints of $e$:
    the main guide with the starting word and the auxiliary guide with the auxiliary word before, and the main guide
    with the first main word and the auxiliary guide with the exchanged
    auxiliary word after. Take their exact finite true-vertex sets and
    an adaptive patch supply at cell centers for $\Omega$, with parameters
    $B_0,C,c$, test exponent $60$ and size $L$. Suppose
    the supplied patches have radius $\epsilon_0t$, with
    $0<\epsilon_0t\le B_0\log L$, and require pairwise-disjoint outer
    physical samples within each of the four vertex sets. Choose
    any owner map to the level-$k$ protocol parties and any natural-valued
    rank. The corresponding supplied guide frames with the protocol
    placeholder have a bounded pair change with error $4L^{-30}$ and
    participants in any set containing the owner of $C_e$, for every
    natural $r,K,r_s$, including zero.

    For the protocol specialization, assume the family, unit state and
    exponent-$60$ angular-scale data just stated, with
    $q\ge1$, $L'_0\le L$, $m_0\ge2$ and $A_0\ge1$. Assume the protocol
    constants $K_0,\epsilon_0,\nu,a,D,B_0$, $L\ge2$, $2^k<2L$, the
    exponent-$60$ adaptive supply, and any label rank. Let $R$ be any
    repaintable square of level $j<k$, index $S_0$ and side
    $n=2^j>K_0t$, where $t=D\log L$. For any set $\mathcal D$ of block
    indices, form the block baseline at $S_0$ from the level-$j$ protocol
    labels at $\mathcal D$, with final label the new square party of $R$.
    For any edge list $E$ and edge $e$, use its neighboring label $C_e$
    to treat all four guides above at the two endpoints of $e$. Their
    treated frames have the same pair-change conclusion for every natural
    $r,K,r_s$, with participants in any set containing the party of $C_e$. The list
    $E$ may repeat edges and may contain $e$; no condition
    $S_0\notin\mathcal D$ is required
    (\texttt{06-geometry.tex}, lines 183--221, 389--430 and 440--445).

\end{theorem}

\begin{proof}\leanok
    \uses{def:peps_frame_exchange_data,
        def:peps_protocol_pair_change,
        lem:peps_adaptive_patch_supply,
        lem:peps_circuit_source_bound_monomials,
        lem:peps_dyadic_floor,
        lem:peps_dyadic_frozen_hole_count,
        lem:peps_dyadic_lens_sides,
        lem:peps_dyadic_schedule_guides,
        thm:peps_bounded_change_closure,
        thm:peps_bulk_steps_treated_frames,
        thm:peps_dyadic_angular_separation,
        thm:peps_dyadic_no_changed_vertex,
        thm:peps_exchange_holes_perm,
        thm:peps_guide_frame_norm,
        thm:peps_guide_frame_supplied_gates,
        thm:peps_pair_change_reorder,
        thm:peps_reading_constant_open,
        thm:peps_reading_piecewise,
        thm:peps_repainted_square_side_bound,
        thm:peps_treated_frame_spec,
        thm:peps_treated_radius_bounds}
    Let $Y$ be the local lens and $Y'=d+Y$ its translate to $R$. Both treated before-guides equal $C_e$ on the open $a_0t$-neighborhood of $\partial Y$, so the after-readings are pieced along $Y'$ (Theorems~\ref{thm:peps_reading_piecewise} and~\ref{thm:peps_reading_constant_open}), every outer hole $\overline B_\infty(z,2\epsilon_0t)$ lies wholly in $Y'$ or $(Y')^c$ (Lemma~\ref{lem:peps_dyadic_lens_sides}), and $Y'$ has diameter $\le n\le L$ (Theorem~\ref{thm:peps_repainted_square_side_bound}). Applying the bulk exchange of Theorem~\ref{thm:peps_bulk_steps_treated_frames} on $Y'$ and reordering the permuted hole lists via Theorems~\ref{thm:peps_exchange_holes_perm} and~\ref{thm:peps_pair_change_reorder} yields the pair change with $0$ nonprivate gates and error $4L^{-30}$.
\end{proof}

\begin{theorem}[Avoiding a value by a small perturbation]
    \label{thm:peps_protocol_affine_avoidance}
    \lean{TNLean.PEPS.Approximation.eventually_add_ne,
        TNLean.PEPS.Approximation.eventually_affine_ne}
    \leanok
    For any $a,m,c\in\R$ with $m\ne0$ (in particular $m=1$), $a+m\eta\ne c$ for all sufficiently small $\eta>0$.
\end{theorem}

\begin{proof}\leanok
    The equation $a+m\eta=c$ has the unique root $\eta=(c-a)/m$, which $(0,\lvert c-a\rvert/\lvert m\rvert)$ (or $(0,\infty)$ when $c=a$) avoids.
\end{proof}

\begin{theorem}[Diagonal avoidance of square boundaries]
    \label{thm:peps_protocol_square_boundary_avoidance}
    \lean{TNLean.PEPS.Approximation.eventually_mem_openSquare_or_notMem_closedSquare}
    \leanok
    \uses{def:peps_dyadic_schedule}
    For any $n\in\R$ and $p\in\R^2$, $p+\eta(1,1)\in(0,n)^2\cup([0,n]^2)^c$ for all sufficiently small $\eta>0$.
\end{theorem}

\begin{proof}\leanok
    \uses{thm:peps_protocol_affine_avoidance, def:peps_dyadic_schedule}
    By Theorem~\ref{thm:peps_protocol_affine_avoidance}, both coordinates of $p+\eta(1,1)$ avoid $\{0,n\}$ for small $\eta>0$, so membership in $[0,n]^2$ forces membership in $(0,n)^2$.
\end{proof}

\begin{theorem}[Diagonal avoidance of band interfaces]
    \label{thm:peps_protocol_band_interface_avoidance}
    \lean{TNLean.PEPS.Approximation.RepaintingBaseline.eventually_notMem_edgeInterface}
    \leanok
    \uses{def:peps_dyadic_schedule, def:peps_dyadic_edge_band}
    For any repainting baseline, edge, and $p\in\R^2$, $p+\eta(1,1)$ avoids that edge's interface for all sufficiently small $\eta>0$ (\texttt{06-geometry.tex}, lines 69--80 and 161--168).
\end{theorem}

\begin{proof}\leanok
    \uses{thm:peps_protocol_affine_avoidance, def:peps_dyadic_edge_band,
        def:peps_dyadic_schedule}
    Along $p+\eta(1,1)$, the normal and tangential coordinates are $a+\sigma\eta$ and $b+\eta$ ($\sigma\in\{\pm1\}$). Since $\sigma,\sigma\mp 1/1000\ne0$, Theorem~\ref{thm:peps_protocol_affine_avoidance} gives $a+\sigma\eta\notin\{0,\min(b+\eta,n-b-\eta)/1000\}$ for small $\eta>0$.
\end{proof}

\begin{theorem}[Reading guides that agree along a displacement]
    \label{thm:peps_protocol_reading_congruence}
    \lean{TNLean.PEPS.Approximation.IsDisplacedReading.congr_eventually}
    \leanok
    \uses{lem:peps_dyadic_sampled_guides}
    If $g$ reads $G$ along $v$ and $G(p+\eta v)=G'(p+\eta v)$ for each $p$ and all sufficiently small $\eta>0$, then $g$ also reads $G'$ along $v$.
\end{theorem}

\begin{proof}\leanok
    \uses{lem:peps_dyadic_sampled_guides}
    Intersecting the two eventual equalities gives $g(p)=G(p+\eta v)=G'(p+\eta v)$ for small $\eta>0$.
\end{proof}

\begin{theorem}[Coordinates within the indexed block]
    \label{thm:peps_protocol_block_coordinates}
    \lean{TNLean.PEPS.Approximation.sub_blockCorner_mem_closedSquare}
    \leanok
    \uses{def:peps_dyadic_block_guide, def:peps_dyadic_schedule}
    If $n>0$ and $x\in\R^2$ has scale-$n$ block index $S=\lfloor x/n\rfloor\in\Z^2$, then $x-nS\in[0,n]^2$.
\end{theorem}

\begin{proof}\leanok
    \uses{def:peps_dyadic_block_guide, def:peps_dyadic_schedule}
    Since $S_i=\lfloor x_i/n\rfloor$, we have $nS_i\le x_i<n(S_i+1)$, so $0\le x_i-nS_i\le n$.
\end{proof}

\begin{definition}[The baseline of a protocol repainting]
    \label{def:peps_protocol_repaint_baseline}
    \lean{TNLean.PEPS.Approximation.repaintBaseline}
    \leanok
    \uses{def:peps_dyadic_protocol_parties, def:peps_dyadic_block_guide,
        def:peps_dyadic_schedule, thm:peps_dyadic_repainting_valid_labels,
        def:peps_dyadic_used_parties}
    For a non-root repaintable square $R$ of level $j<k$, index $S$, and side $n=2^j$ at stage $\mathcal D\subseteq\Z^2$, its repainting baseline $\mathcal B$ is the block baseline at $S$ with the level-$j$ labels at $\mathcal D$ and the fresh label of $R$; its starting guide $g_{\mathrm{start}}(x)$ is the level-$j$ block guide evaluated at $x+nS$ (\texttt{06-geometry.tex}, lines 82--125 and 142--160).
\end{definition}

\begin{theorem}[The guides of an unmarked repainting]
    \label{thm:peps_protocol_unmarked_guides}
    \lean{TNLean.PEPS.Approximation.treatedGuide_completedGuide,
        TNLean.PEPS.Approximation.treatedGuide_empty}
    \leanok
    \uses{def:peps_protocol_repaint_baseline, def:peps_protocol_treated_frame,
        def:peps_protocol_level_frame, def:peps_dyadic_block_guide,
        def:peps_dyadic_schedule}
    Treating any guide at $M=\varnothing$ (Definition~\ref{def:peps_protocol_treated_frame}) yields the sampled translate of the guide to the block $2^iT+[0,2^i]^2$. For the baseline $\mathcal B$ of $R$ at $\mathcal D$ and any list $E$ containing all four edges, the completed guide $g_{\mathrm{done}}(E)$ placed on $R$ and treated at $M=\varnothing$ equals the level-$j$ guide at $\mathcal D\cup\{S\}$ (\texttt{06-geometry.tex}, lines 69--80 and 215--221).
\end{theorem}

\begin{proof}\leanok
    \uses{def:peps_protocol_treated_frame, def:peps_protocol_level_frame,
        def:peps_protocol_repaint_baseline, def:peps_dyadic_block_guide,
        lem:peps_dyadic_floor, thm:peps_protocol_reading_congruence,
        thm:peps_protocol_square_boundary_avoidance,
        thm:peps_protocol_block_coordinates, lem:peps_dyadic_schedule_guides,
        lem:peps_dyadic_active_labels, thm:peps_dyadic_repainting_valid_labels,
        thm:peps_protocol_level_frame_facts,
        thm:peps_protocol_bounded_changes}
    At $M=\varnothing$ the homogenized region is empty. Once all four edges are in $E$, $g_{\mathrm{done}}(E)$ takes the fresh label of $R$ on $(0,n)^2$ and the starting label on $([0,n]^2)^c$ (Lemma~\ref{lem:peps_dyadic_schedule_guides}); Theorems~\ref{thm:peps_protocol_square_boundary_avoidance}, \ref{thm:peps_protocol_block_coordinates}, and~\ref{thm:peps_protocol_reading_congruence} identify its sampled translate with the level-$j$ guide at $\mathcal D\cup\{S\}$.
\end{proof}

\begin{theorem}[Treated frames of equal treated guides]
    \label{thm:peps_protocol_treated_frame_congruence}
    \lean{TNLean.PEPS.Approximation.treatedFrame_congr,
        TNLean.PEPS.Approximation.treatedFrame_empty_congr}
    \leanok
    \uses{def:peps_protocol_treated_frame, def:peps_adaptive_patch_constraint,
        def:peps_dyadic_protocol_parties}
    In Definition~\ref{def:peps_protocol_treated_frame}, for $0<\epsilon_0t\le B_0\log L$, equal treated guides with the same finite mark set $M$ produce equal treated frames $F_M^\lambda(g)$; in particular, $F_\varnothing^\lambda(g)$ is independent of the homogenizing label $\lambda$.
\end{theorem}

\begin{proof}\leanok
    \uses{def:peps_protocol_treated_frame, thm:peps_protocol_unmarked_guides}
    By Definition~\ref{def:peps_protocol_treated_frame}, $F_M^\lambda(g)$ depends on $(g,\lambda)$ only through the treated guide, which at $M=\varnothing$ is independent of $\lambda$ by Theorem~\ref{thm:peps_protocol_unmarked_guides}.
\end{proof}

\begin{theorem}[The treated frames at the ends of a large repainting]
    \label{thm:peps_protocol_treated_frame_ends}
    \lean{TNLean.PEPS.Approximation.levelFrame_eq_treatedFrame_guide,
        TNLean.PEPS.Approximation.treatedFrame_completedGuide_eq_levelFrame,
        TNLean.PEPS.Approximation.treatedFrame_edgeStart,
        TNLean.PEPS.Approximation.treatedFrame_auxStart_eq_raw}
    \leanok
    \uses{def:peps_protocol_repaint_baseline, def:peps_protocol_treated_frame,
        def:peps_protocol_level_frame, def:peps_adaptive_patch_constraint,
        def:peps_dyadic_protocol_parties, def:peps_dyadic_schedule,
        def:peps_dyadic_edge_band, thm:peps_protocol_constants, def:peps_frame,
        def:peps_dyadic_normal_words,
        thm:peps_treated_radius_bounds}
    Let $R$ be a non-root repaintable square of level $j<k$, index $S$, side $n=2^j$, and baseline $\mathcal B$ at stage $\mathcal D\subseteq\Z^2$ (Definition~\ref{def:peps_protocol_repaint_baseline}).
    \begin{enumerate}
        \item For $0<\epsilon_0t\le B_0\log L$ and any edge $e$ with neighbor label $C_e$, the unmarked treated frame of the uniform auxiliary guide at $C_e$ is the hole-free raw copy owned by the party of $C_e$.
        \item Under protocol constants with $2\le L$, $t=D\log L$, and $n>K_0t$, for any edge list $E$, $e\notin E$, and labels $\lambda_1,\lambda_2$, the unmarked frame $F_\varnothing^{\lambda_1}(g_{\mathrm{done}}(E))$ equals $F_\varnothing^{\lambda_2}(g_{\mathrm{main}}(E,e,w_{\mathrm{start}}))$, where $w_{\mathrm{start}}$ is the starting band word along $e$.
        \item If also $L\le 2^k$, then for any non-root reference square $R_0$, the level-$j$ main frame at $\mathcal D$ (with hole radius $h_n=\epsilon_0t$ and party map of $R_0$) equals $F_\varnothing^\lambda(g_{\mathrm{start}})$, and for any list $F$ containing all four edges, $F_\varnothing^\lambda(g_{\mathrm{done}}(F))$ equals the level-$j$ main frame at $\mathcal D\cup\{S\}$ (\texttt{06-geometry.tex}, lines 62--125, 142--221 and 328--339).
    \end{enumerate}
\end{theorem}

\begin{proof}\leanok
    \uses{thm:peps_protocol_unmarked_guides, def:peps_protocol_repaint_baseline,
        thm:peps_protocol_treated_frame_congruence,
        thm:peps_protocol_reading_congruence,
        thm:peps_protocol_band_interface_avoidance,
        thm:peps_protocol_party_map_switch, thm:peps_protocol_level_frame_facts,
        lem:peps_dyadic_schedule_guides, thm:peps_guide_frame_supplied_gates,
        lem:peps_dyadic_floor, lem:peps_dyadic_sampled_guides,
        lem:peps_dyadic_incident_change, def:peps_protocol_treated_frame,
        thm:peps_protocol_level_events, thm:peps_protocol_constants,
        def:peps_dyadic_block_guide, def:peps_dyadic_edge_band,
        def:peps_dyadic_schedule, def:peps_dyadic_normal_words,
        def:peps_guide_frame, def:peps_frame, def:peps_protocol_level_frame,
        thm:peps_treated_radius_bounds, thm:peps_treated_frame_spec}
    The uniform guide at $C_e$ has no true vertices, giving (1). For (2), $g_{\mathrm{done}}(E)$ and $g_{\mathrm{main}}(E,e,w_{\mathrm{start}})$ agree off the interface of $e$ (Lemma~\ref{lem:peps_dyadic_schedule_guides}), so Theorems~\ref{thm:peps_protocol_band_interface_avoidance}, \ref{thm:peps_protocol_reading_congruence}, and~\ref{thm:peps_protocol_treated_frame_congruence} identify their unmarked frames. For (3), Theorem~\ref{thm:peps_protocol_unmarked_guides} identifies the sampled guides with the level-$j$ guides at $\mathcal D$ and $\mathcal D\cup\{S\}$, $h_n=\epsilon_0\min(n,t)=\epsilon_0t$ matches the patch radii, and Theorem~\ref{thm:peps_protocol_party_map_switch} switches the party map to $R_0$.
\end{proof}

\begin{definition}[A chain of a prescribed number of steps]
    \label{def:peps_protocol_step_chain}
    \lean{TNLean.PEPS.Approximation.IsStepChain}
    \leanok
    For a binary relation $\mathcal R$ on $X$ and $n\in\mathbb N$, a chain of $n$ steps from $a\in X$ to $b\in X$ is the existence of $s\colon\mathbb N\to X$ with
    \begin{align}
        s_0 & = a,\qquad s_n=b,\qquad
        \forall j<n,\quad \mathcal R(s_j,s_{j+1}). \notag
    \end{align}
\end{definition}

\begin{lemma}[Chains of changes]
    \label{lem:peps_protocol_change_chains}
    \lean{TNLean.PEPS.Approximation.IsStepChain.refl,
        TNLean.PEPS.Approximation.IsStepChain.single,
        TNLean.PEPS.Approximation.IsStepChain.trans,
        TNLean.PEPS.Approximation.IsStepChain.map,
        TNLean.PEPS.Approximation.IsStepChain.cast_length,
        TNLean.PEPS.Approximation.IsStepChain.treat,
        TNLean.PEPS.Approximation.IsStepChain.untreat,
        TNLean.PEPS.Approximation.IsStepChain.treat_operate,
        TNLean.PEPS.Approximation.IsStepChain.treatOperate,
        TNLean.PEPS.Approximation.isStepChain_main,
        TNLean.PEPS.Approximation.hasSourceBoundedChange_aux_of_eq,
        TNLean.PEPS.Approximation.isStepChain_aux,
        TNLean.PEPS.Approximation.hasSourceBoundedChange_of_hasPairChange}
    \leanok
    \uses{def:peps_protocol_step_chain, def:peps_frame_change,
        thm:peps_protocol_event_segments, def:peps_protocol_pair_change}
    Step chains (Definition~\ref{def:peps_protocol_step_chain}) contain $0$-step reflexive chains and $1$-step single transitions, compose by adding lengths ($m+n$), and push forward under relation-preserving maps. Given mark-indexed states $T(M),T'(M)$ and a mark list of length $m$ with underlying set $M_*$, if adding (for $T$) and removing (for $T'$) each listed mark is a step and $\mathcal R(T(M_*),T'(M_*))$, then $T(\varnothing)$ reaches $T'(\varnothing)$ in $2m+1$ steps. When $q\ge1$ and $\lVert\Omega\rVert\le1$, chains of bounded changes (Definitions~\ref{def:peps_frame_change} and~\ref{def:peps_protocol_pair_change}) of the main sheet, of an auxiliary sheet in a duplicate-free index list, or of both together lift to the full sheet memory with the same bounds $(S,r,K,r_s,\varepsilon)$ (\texttt{06-geometry.tex}, lines 328--339 and 572--583).
\end{lemma}

\begin{proof}\leanok
    \uses{def:peps_protocol_step_chain, thm:peps_protocol_level_frame_facts,
        def:peps_frame_change, def:peps_protocol_pair_change}
    Induction along the mark list gives $m$-step treatment and untreatment chains, which concatenate with the middle step to length $m+1+m=2m+1$; the memory lifts follow from Theorem~\ref{thm:peps_protocol_level_frame_facts}.
\end{proof}

\begin{lemma}[Enlarging the error bound of frame changes]
    \label{lem:peps_protocol_frame_change_error_mono}
    \lean{TNLean.PEPS.Approximation.HasFrameChange.mono_error,
        TNLean.PEPS.Approximation.HasPairChange.mono_error}
    \leanok
    \uses{def:peps_frame,
        def:peps_frame_change,
        def:peps_protocol_pair_change}
    For $\varepsilon\le\varepsilon'$, any bounded frame or pair change (Definitions~\ref{def:peps_frame_change} and~\ref{def:peps_protocol_pair_change}) with bounds $(S,r,K,r_s,\varepsilon)$ is also one with bounds $(S,r,K,r_s,\varepsilon')$.
\end{lemma}

\begin{proof}\leanok
    \uses{def:peps_frame_change,
        def:peps_protocol_pair_change,
        thm:peps_bounded_change_closure}
    Immediate from the monotonicity of bounded changes (Theorem~\ref{thm:peps_bounded_change_closure}).
\end{proof}

\begin{definition}[The corner and edge-endpoint lists]
    \label{def:peps_protocol_bulk_operation_marks}
    \lean{TNLean.PEPS.Approximation.cornerList,
        TNLean.PEPS.Approximation.SquareEdge.endpointList,
        TNLean.PEPS.Approximation.toFinset_cornerList,
        TNLean.PEPS.Approximation.SquareEdge.toFinset_endpointList,
        TNLean.PEPS.Approximation.isBlockCornerOf_of_mem_cornerList,
        TNLean.PEPS.Approximation.SquareEdge.isBlockCornerOf_of_mem_endpointList}
    \leanok
    \uses{def:peps_dyadic_edge_band,
        lem:peps_dyadic_direct_repainting,
        thm:peps_bulk_steps_treated_frames}
    For a block index $S\in\mathbb Z^2$, its ordered corner list is $[S,S+(1,0),S+(0,1),S+(1,1)]$, and its bottom, right, top, and left edge endpoint lists are respectively
    \begin{align}
         & [S,S+(1,0)],\quad [S+(1,0),S+(1,1)],\quad
        [S+(0,1),S+(1,1)],\quad [S,S+(0,1)], \notag
    \end{align}
    whose underlying finite sets are the $4$-corner set of $S$ and the $2$-endpoint sets of each edge.
\end{definition}

\begin{proof}\leanok
    \uses{def:peps_dyadic_edge_band,
        lem:peps_dyadic_direct_repainting,
        thm:peps_bulk_steps_treated_frames}
    Immediate by inspection of the four corners and edges.
\end{proof}

\begin{definition}[Local labels of a repainting]
    \label{def:peps_protocol_local_labels}
    \lean{TNLean.PEPS.Approximation.IsLocalLabel,
        TNLean.PEPS.Approximation.isHomogenizing_oldLabel,
        TNLean.PEPS.Approximation.isHomogenizing_finalLabel,
        TNLean.PEPS.Approximation.isHomogenizing_nbrLabel}
    \leanok
    \uses{def:peps_dyadic_edge_band,
        def:peps_dyadic_protocol_parties,
        def:peps_dyadic_used_parties,
        def:peps_protocol_repaint_baseline,
        thm:peps_dyadic_block_ratio,
        thm:peps_dyadic_repainting_valid_labels,
        thm:peps_protocol_bounded_changes}
    For a repaintable square $R$ of index $S$ at stage $\mathcal D\subseteq\mathbb Z^2$, a guide label is local if it is the new label of $R$ or the level label at $\mathcal D$ of a block in $S+\{-1,0,1\}^2$. In particular, the old label $A$, the final label $B$, and the neighbor label $C_e$ across each edge $e$ of $R$ are local.
\end{definition}

\begin{proof}\leanok
    \uses{def:peps_dyadic_block_guide,
        def:peps_dyadic_edge_band,
        def:peps_protocol_repaint_baseline,
        thm:peps_dyadic_point_ratio,
        thm:peps_dyadic_repainting_valid_labels,
        thm:peps_protocol_bounded_changes}
    The old label corresponds to offset $(0,0)\in\{-1,0,1\}^2$, the final label is the new label of $R$, and each edge neighbor lies in $S+\{-1,0,1\}^2$.
\end{proof}

\begin{definition}[Bounded point treatments]
    \label{def:peps_protocol_local_point_treatments}
    \lean{TNLean.PEPS.Approximation.IsPointTreatable}
    \leanok
    \uses{def:peps_adaptive_patch_constraint,
        def:peps_dyadic_protocol_parties,
        def:peps_frame_change,
        def:peps_protocol_local_labels,
        def:peps_protocol_repaint_baseline,
        def:peps_protocol_treated_frame,
        def:polynomial_peps_model,
        lem:peps_adaptive_patch_supply,
        lem:peps_dyadic_direct_repainting,
        lem:peps_dyadic_unmodified_vertices,
        thm:peps_protocol_bounded_changes}
    Given $q\ge1$, an adaptive patch supply $P$ for $\Omega$ with parameters $(B_0,C,c,60,L)$, a label rank, a repaintable square $R$ at stage $\mathcal D\subseteq\mathbb Z^2$, and $0<\epsilon_0t\le B_0\log L$, the repainting has bounded point treatments with participants $S$ and bounds $(r,K,r_s,\varepsilon)$ if for every unmodified baseline guide $g$, corner $Q_0$ of $R$, finite mark sets $M,M'\subseteq\mathbb Z^2$ agreeing away from $Q_0$, and local label $\lambda$ (Definition~\ref{def:peps_protocol_local_labels}), the treated frames $F_M^\lambda(g)$ and $F_{M'}^\lambda(g)$ (Definition~\ref{def:peps_protocol_treated_frame}) are joined by a bounded frame change with parameters $(S,r,K,r_s,\varepsilon)$.
\end{definition}

\begin{theorem}[Comparison with the small-patch error]
    \label{thm:peps_protocol_small_patch_error_bounds}
    \lean{TNLean.PEPS.Approximation.zpow_neg_thirty_le,
        TNLean.PEPS.Approximation.four_mul_zpow_neg_thirty_le}
    \leanok
    For any integer $L\ge1$, $L^{-30}\le(2b+1)L^{-20}$ for all $b\in\mathbb N$, and $4L^{-30}\le(2b+1)L^{-20}$ for all $b\ge2$.
\end{theorem}

\begin{proof}\leanok
    Since $L\ge1$, we have $0\le L^{-30}\le L^{-20}$, and $2b+1\ge1$ for $b\ge0$ while $2b+1\ge5\ge4$ for $b\ge2$.
\end{proof}

\begin{theorem}[A sheet operation with an arbitrary list of treatments]
    \label{thm:peps_protocol_sheet_operation_chains}
    \lean{TNLean.PEPS.Approximation.isStepChain_sheetOperation}
    \leanok
    \uses{def:peps_adaptive_patch_constraint,
        def:peps_dyadic_protocol_parties,
        def:peps_frame_change,
        def:peps_protocol_local_labels,
        def:peps_protocol_local_point_treatments,
        def:peps_protocol_repaint_baseline,
        def:peps_protocol_step_chain,
        def:peps_protocol_treated_frame,
        def:polynomial_peps_model,
        lem:peps_adaptive_patch_supply,
        lem:peps_dyadic_direct_repainting,
        lem:peps_dyadic_unmodified_vertices,
        thm:peps_protocol_bounded_changes}
    Assume bounded point treatments with parameters $(S,r,K,r_s,\varepsilon)$ as in Definition~\ref{def:peps_protocol_local_point_treatments}. For any unmodified guides $g,g'$, local label $\lambda$, and list of $m$ corners of $R$ with underlying finite set $M$, a bounded frame change from $F_M^\lambda(g)$ to $F_M^\lambda(g')$ with bounds $(S,r,K,r_s,\varepsilon)$ extends to a chain of $2m+1$ such bounded frame changes from $F_\varnothing^\lambda(g)$ to $F_\varnothing^\lambda(g')$.
\end{theorem}

\begin{proof}\leanok
    \uses{def:peps_protocol_local_point_treatments,
        def:peps_protocol_treated_frame,
        lem:peps_protocol_change_chains}
    Adding or removing one listed corner changes mark membership only at that corner, which Definition~\ref{def:peps_protocol_local_point_treatments} implements; Lemma~\ref{lem:peps_protocol_change_chains} combines the $m$ treatments, the bulk step on $M$, and the $m$ untreatments into $2m+1$ steps.
\end{proof}

\begin{theorem}[The party of a local label participates]
    \label{thm:peps_protocol_local_label_event_uses}
    \lean{TNLean.PEPS.Approximation.IsLocalLabel.eventUses}
    \leanok
    \uses{def:peps_dyadic_protocol_parties,
        def:peps_protocol_local_labels}
    For any repaintable square $R$ at stage $\mathcal D\subseteq\mathbb Z^2$, the party of every local label $\lambda$ (Definition~\ref{def:peps_protocol_local_labels}) belongs to the participant set $S_R$ of the block event $R$.
\end{theorem}

\begin{proof}\leanok
    \uses{def:peps_protocol_local_labels,
        thm:peps_protocol_bounded_changes}
    Immediate from Definition~\ref{def:peps_protocol_local_labels} and Theorem~\ref{thm:peps_protocol_bounded_changes}.
\end{proof}

\begin{theorem}[The central birth with its point treatments]
    \label{thm:peps_protocol_central_birth_chain}
    \lean{TNLean.PEPS.Approximation.isStepChain_centralBirth}
    \leanok
    \uses{def:peps_adaptive_patch_constraint,
        def:peps_angular_scale,
        def:peps_dyadic_protocol_parties,
        def:peps_dyadic_schedule,
        def:peps_frame_change,
        def:peps_protocol_local_point_treatments,
        def:peps_protocol_repaint_baseline,
        def:peps_protocol_step_chain,
        def:peps_protocol_treated_frame,
        def:polynomial_peps_model,
        lem:peps_adaptive_patch_supply,
        thm:peps_protocol_constants}
    Let $k,m_0,L,b,K,q\in\mathbb N$ with $q\ge1$, and let all the scale
    constants below be real. Fix a family of states and assume the angular-scale
    condition for it with parameters $q,a,A_0,m_0,60,D,L_*$ and $L_*\le L$. Let
    $\Omega$ be a unit vector of the family on the square of side $L$. Assume $m_0\ge4$,
    $A_0\ge1$, the protocol-constant condition on $K_0,\epsilon_0,\nu,a,D,B_0$,
    $L\ge2$ and $2^k<2L$. Set $t=D\log L$. Fix an adaptive patch supply
    $P$ at the cell-center positions for $\Omega$ with parameters
    $B_0,C,c,60,L$, and an arbitrary natural-valued label rank.
    Let $R$ be a repaintable square of hierarchy $k$, with side $n=2^j>K_0t$,
    and let $\mathit{done}\subseteq\mathbb Z^2$ be arbitrary. Assume
    $0<\epsilon_0t\le B_0\log L$, so these treated frames use supplied
    patches of radius $\epsilon_0t$.
    Let $S$ be the participant set of the block event $R$, put
    $\delta_b=(2b+1)(L^{20})^{-1}$, and assume $b\ge1$ and $K\ge1$.
    Assume bounded point treatments for these data with participants $S$ and
    bounds $(2b,K,2b,\delta_b)$. In the following conclusion each step has
    these same bounds. Write $\mathcal B$ for the repainting baseline,
    $A,B$ for its old and final labels, and $F_M^\lambda(g)$ for its treated
    frame at marks $M$ with homogenizing label $\lambda$.
    Then there is a chain of $9$ bounded frame changes from
    $F_\varnothing^A(\mathcal B.\mathrm{guide})$ to
    $F_\varnothing^A(\mathcal B.\mathrm{centralGuide})$. It consists of the
    four corner treatments with old label $A$, the central birth, and removal
    of the four treatments.
\end{theorem}

\begin{proof}\leanok
    \uses{def:peps_dyadic_protocol_parties,
        def:peps_protocol_bulk_operation_marks,
        def:peps_protocol_local_labels,
        lem:peps_protocol_frame_change_error_mono,
        thm:peps_protocol_bounded_changes,
        thm:peps_protocol_central_birth_change,
        thm:peps_protocol_local_label_event_uses,
        thm:peps_protocol_sheet_operation_chains,
        thm:peps_protocol_small_patch_error_bounds}
    By Definition~\ref{def:peps_protocol_local_labels} and Theorem~\ref{thm:peps_protocol_local_label_event_uses}, the parties of $A$ and $B$ belong to $S$. Enlarging the central-birth error $L^{-30}\le\delta_b$ via Theorem~\ref{thm:peps_protocol_small_patch_error_bounds} and Lemma~\ref{lem:peps_protocol_frame_change_error_mono}, Theorem~\ref{thm:peps_protocol_sheet_operation_chains} applied to the $4$-corner list of Definition~\ref{def:peps_protocol_bulk_operation_marks} yields the $9$-step chain.
\end{proof}

\begin{theorem}[An elementary band operation with its point treatments]
    \label{thm:peps_protocol_band_operation_chain}
    \lean{TNLean.PEPS.Approximation.isStepChain_bandOperation}
    \leanok
    \uses{def:peps_adaptive_patch_constraint,
        def:peps_angular_scale,
        def:peps_dyadic_edge_band,
        def:peps_dyadic_normal_words,
        def:peps_dyadic_protocol_parties,
        def:peps_dyadic_schedule,
        def:peps_frame_change,
        def:peps_protocol_local_labels,
        def:peps_protocol_local_point_treatments,
        def:peps_protocol_repaint_baseline,
        def:peps_protocol_step_chain,
        def:peps_protocol_treated_frame,
        def:polynomial_peps_model,
        lem:peps_adaptive_patch_supply,
        lem:peps_dyadic_unmodified_vertices,
        thm:peps_dyadic_no_changed_vertex,
        thm:peps_protocol_constants}
    Let $k,m_0,L,b,K,q\in\mathbb N$ with $q\ge1$, and let all the scale
    constants below be real. Fix a family of states and assume the angular-scale
    condition for it with parameters $q,a,A_0,m_0,60,D,L_*$ and $L_*\le L$. Let
    $\Omega$ be a unit vector of the family on the square of side $L$. Assume $m_0\ge2$,
    $A_0\ge1$, the protocol-constant condition on $K_0,\epsilon_0,\nu,a,D,B_0$,
    $L\ge2$ and $2^k<2L$. Set $t=D\log L$. Fix an adaptive patch supply
    $P$ at the cell-center positions for $\Omega$ with parameters
    $B_0,C,c,60,L$, and an arbitrary natural-valued label rank.
    Let $R$ be a repaintable square of hierarchy $k$, with side $n=2^j>K_0t$,
    and let $\mathit{done}\subseteq\mathbb Z^2$ be arbitrary. Assume
    $0<\epsilon_0t\le B_0\log L$, so these treated frames use supplied
    patches of radius $\epsilon_0t$.
    Let $S$ be the participant set of the block event $R$, put
    $\delta_b=(2b+1)(L^{20})^{-1}$, and assume $b\ge1$ and $K\ge1$.
    Assume bounded point treatments for these data with participants $S$ and
    bounds $(2b,K,2b,\delta_b)$. In the following conclusion each step has
    these same bounds. Write $\mathcal B$ for the repainting baseline,
    $A,B$ for its old and final labels, and $F_M^\lambda(g)$ for its treated
    frame at marks $M$ with homogenizing label $\lambda$.
    Let $\mathit{op}$ be an elementary band operation that inserts or removes
    one label $Q$, in the sense of the one-label operation predicate. Assume
    both $Q$ and the surrounding label $\mathit{op}.\mathrm{label}$ are local.
    For any edge $e$ and arbitrary background guide $g_0$, suppose both
    $g_-=\mathrm{bandUpdate}(n,e,\mathit{op}.\mathrm{before},g_0)$ and
    $g_+=\mathrm{bandUpdate}(n,e,\mathit{op}.\mathrm{after},g_0)$ are unmodified
    guides of $\mathcal B$. Then there is a chain of $5$ bounded frame changes
    from $F_\varnothing^{\mathit{op}.\mathrm{label}}(g_-)$ to
    $F_\varnothing^{\mathit{op}.\mathrm{label}}(g_+)$: the two endpoint
    treatments, the operation, and removal of the treatments.
\end{theorem}

\begin{proof}\leanok
    \uses{def:peps_dyadic_edge_band,
        def:peps_dyadic_protocol_parties,
        def:peps_dyadic_schedule,
        def:peps_frame_change,
        def:peps_protocol_bulk_operation_marks,
        def:peps_protocol_treated_frame,
        lem:peps_protocol_change_chains,
        lem:peps_protocol_frame_change_error_mono,
        thm:peps_bulk_steps_treated_frames,
        thm:peps_protocol_band_death_change,
        thm:peps_protocol_band_step_changes,
        thm:peps_protocol_bounded_changes,
        thm:peps_protocol_local_label_event_uses,
        thm:peps_protocol_sheet_operation_chains,
        thm:peps_protocol_small_patch_error_bounds}
    Theorem~\ref{thm:peps_protocol_band_step_changes} (birth) or Theorem~\ref{thm:peps_protocol_band_death_change} (death) gives the endpoint-treated frame change with error $L^{-30}\le\delta_b$ (Theorem~\ref{thm:peps_protocol_small_patch_error_bounds} and Lemma~\ref{lem:peps_protocol_frame_change_error_mono}) and participants in $S$ (Theorem~\ref{thm:peps_protocol_local_label_event_uses}); Theorem~\ref{thm:peps_protocol_sheet_operation_chains} applied to the $2$-endpoint list of $e$ (Definition~\ref{def:peps_protocol_bulk_operation_marks}) gives $2\cdot2+1=5$ steps.
\end{proof}

\begin{theorem}[The three auxiliary births]
    \label{thm:peps_protocol_auxiliary_birth_chain}
    \lean{TNLean.PEPS.Approximation.isStepChain_auxBirths}
    \leanok
    \uses{def:peps_adaptive_patch_constraint,
        def:peps_angular_scale,
        def:peps_dyadic_edge_band,
        def:peps_dyadic_normal_words,
        def:peps_dyadic_protocol_parties,
        def:peps_dyadic_schedule,
        def:peps_frame_change,
        def:peps_protocol_local_point_treatments,
        def:peps_protocol_repaint_baseline,
        def:peps_protocol_step_chain,
        def:peps_protocol_treated_frame,
        def:polynomial_peps_model,
        lem:peps_adaptive_patch_supply,
        thm:peps_protocol_constants}
    Let $k,m_0,L,b,K,q\in\mathbb N$ with $q\ge1$, and let all the scale
    constants below be real. Fix a family of states and assume the angular-scale
    condition for it with parameters $q,a,A_0,m_0,60,D,L_*$ and $L_*\le L$. Let
    $\Omega$ be a unit vector of the family on the square of side $L$. Assume $m_0\ge2$,
    $A_0\ge1$, the protocol-constant condition on $K_0,\epsilon_0,\nu,a,D,B_0$,
    $L\ge2$ and $2^k<2L$. Set $t=D\log L$. Fix an adaptive patch supply
    $P$ at the cell-center positions for $\Omega$ with parameters
    $B_0,C,c,60,L$, and an arbitrary natural-valued label rank.
    Let $R$ be a repaintable square of hierarchy $k$, with side $n=2^j>K_0t$,
    and let $\mathit{done}\subseteq\mathbb Z^2$ be arbitrary. Assume
    $0<\epsilon_0t\le B_0\log L$, so these treated frames use supplied
    patches of radius $\epsilon_0t$.
    Let $S$ be the participant set of the block event $R$, put
    $\delta_b=(2b+1)(L^{20})^{-1}$, and assume $b\ge1$ and $K\ge1$.
    Assume bounded point treatments for these data with participants $S$ and
    bounds $(2b,K,2b,\delta_b)$. In the following conclusion each step has
    these same bounds. Write $\mathcal B$ for the repainting baseline,
    $A,B$ for its old and final labels, and $F_M^\lambda(g)$ for its treated
    frame at marks $M$ with homogenizing label $\lambda$.
    For any edge $e$ and any homogenizing label $\lambda$, put
    $C_e=\mathcal B.\mathrm{nbrLabel}(e)$. There is a chain of $15$ bounded
    frame changes, with both endpoints treated at no mark, from the auxiliary
    guide of $e$ with the empty band word over $C_e$ to its auxiliary guide with
    word $\mathrm{auxWord}(A,B,C_e)$. Both endpoint frames use the arbitrary
    label $\lambda$; it need not be local.
\end{theorem}

\begin{proof}\leanok
    \uses{def:peps_dyadic_schedule,
        def:peps_protocol_local_labels,
        def:peps_protocol_repaint_baseline,
        lem:peps_dyadic_unmodified_vertices,
        lem:peps_protocol_change_chains,
        thm:peps_dyadic_no_changed_vertex,
        thm:peps_protocol_band_operation_chain,
        thm:peps_protocol_treated_frame_congruence}
    The labels $A,B,C_e$ are local (Definition~\ref{def:peps_protocol_local_labels}) and all intermediate auxiliary words produce unmodified guides (Lemma~\ref{lem:peps_dyadic_unmodified_vertices}). Applying Theorem~\ref{thm:peps_protocol_band_operation_chain} to each of the $3$ auxiliary births and switching the unmarked homogenizing label to $\lambda$ via Theorem~\ref{thm:peps_protocol_treated_frame_congruence} gives $3\cdot 5=15$ steps (Lemma~\ref{lem:peps_protocol_change_chains}).
\end{proof}

\begin{theorem}[The four main-sheet operations]
    \label{thm:peps_protocol_main_steps_chain}
    \lean{TNLean.PEPS.Approximation.isStepChain_mainSteps}
    \leanok
    \uses{def:peps_adaptive_patch_constraint,
        def:peps_angular_scale,
        def:peps_dyadic_edge_band,
        def:peps_dyadic_normal_words,
        def:peps_dyadic_protocol_parties,
        def:peps_dyadic_schedule,
        def:peps_frame_change,
        def:peps_protocol_local_point_treatments,
        def:peps_protocol_repaint_baseline,
        def:peps_protocol_step_chain,
        def:peps_protocol_treated_frame,
        def:polynomial_peps_model,
        lem:peps_adaptive_patch_supply,
        thm:peps_protocol_constants}
    Let $k,m_0,L,b,K,q\in\mathbb N$ with $q\ge1$, and let all the scale
    constants below be real. Fix a family of states and assume the angular-scale
    condition for it with parameters $q,a,A_0,m_0,60,D,L_*$ and $L_*\le L$. Let
    $\Omega$ be a unit vector of the family on the square of side $L$. Assume $m_0\ge2$,
    $A_0\ge1$, the protocol-constant condition on $K_0,\epsilon_0,\nu,a,D,B_0$,
    $L\ge2$ and $2^k<2L$. Set $t=D\log L$. Fix an adaptive patch supply
    $P$ at the cell-center positions for $\Omega$ with parameters
    $B_0,C,c,60,L$, and an arbitrary natural-valued label rank.
    Let $R$ be a repaintable square of hierarchy $k$, with side $n=2^j>K_0t$,
    and let $\mathit{done}\subseteq\mathbb Z^2$ be arbitrary. Assume
    $0<\epsilon_0t\le B_0\log L$, so these treated frames use supplied
    patches of radius $\epsilon_0t$.
    Let $S$ be the participant set of the block event $R$, put
    $\delta_b=(2b+1)(L^{20})^{-1}$, and assume $b\ge1$ and $K\ge1$.
    Assume bounded point treatments for these data with participants $S$ and
    bounds $(2b,K,2b,\delta_b)$. In the following conclusion each step has
    these same bounds. Write $\mathcal B$ for the repainting baseline,
    $A,B$ for its old and final labels, and $F_M^\lambda(g)$ for its treated
    frame at marks $M$ with homogenizing label $\lambda$.
    Let $E$ be any list of square edges, $e$ any edge, and $\lambda$ any guide
    label. Put $C_e=\mathcal B.\mathrm{nbrLabel}(e)$. There is a chain of
    $20$ bounded frame changes from
    $F_\varnothing^\lambda(\mathcal B.\mathrm{mainGuide}
            (E,e,\mathrm{mainWordOne}(A,B,C_e)))$ to
    $F_\varnothing^\lambda(\mathcal B.\mathrm{completedGuide}(e::E))$.
    No freshness or distinctness condition is imposed on $E$ or $e$, and
    $\lambda$ need not be local.
\end{theorem}

\begin{proof}\leanok
    \uses{def:peps_dyadic_schedule,
        def:peps_protocol_local_labels,
        def:peps_protocol_repaint_baseline,
        lem:peps_dyadic_unmodified_vertices,
        lem:peps_protocol_change_chains,
        thm:peps_dyadic_no_changed_vertex,
        thm:peps_protocol_band_operation_chain,
        thm:peps_protocol_treated_frame_congruence}
    Applying Theorem~\ref{thm:peps_protocol_band_operation_chain} to the $4$ elementary band operations on the main sheet (first death, second death, birth, final death) and replacing the unmarked homogenizing labels by $\lambda$ via Theorem~\ref{thm:peps_protocol_treated_frame_congruence} yields $4\cdot 5=20$ steps (Lemma~\ref{lem:peps_protocol_change_chains}).
\end{proof}

\begin{theorem}[The lens exchange with both sheets treated]
    \label{thm:peps_protocol_lens_exchange_chain}
    \lean{TNLean.PEPS.Approximation.isStepChain_lensExchange}
    \leanok
    \uses{def:peps_adaptive_patch_constraint,
        def:peps_angular_scale,
        def:peps_dyadic_edge_band,
        def:peps_dyadic_normal_words,
        def:peps_dyadic_protocol_parties,
        def:peps_dyadic_schedule,
        def:peps_frame,
        def:peps_protocol_local_point_treatments,
        def:peps_protocol_repaint_baseline,
        def:peps_protocol_step_chain,
        def:peps_protocol_treated_frame,
        def:polynomial_peps_model,
        lem:peps_adaptive_patch_supply,
        lem:peps_dyadic_unmodified_vertices,
        thm:peps_protocol_constants,
        thm:peps_protocol_event_segments}
    Let $k,m_0,L,b,K,q\in\mathbb N$ with $q\ge1$, and let all the scale
    constants below be real. Fix a family of states and assume the angular-scale
    condition for it with parameters $q,a,A_0,m_0,60,D,L_*$ and $L_*\le L$. Let
    $\Omega$ be a unit vector of the family on the square of side $L$. Assume $m_0\ge2$,
    $A_0\ge1$, the protocol-constant condition on $K_0,\epsilon_0,\nu,a,D,B_0$,
    $L\ge2$ and $2^k<2L$. Set $t=D\log L$. Fix an adaptive patch supply
    $P$ at the cell-center positions for $\Omega$ with parameters
    $B_0,C,c,60,L$, and an arbitrary natural-valued label rank.
    Let $R$ be a repaintable square of hierarchy $k$, with side $n=2^j>K_0t$,
    and let $\mathit{done}\subseteq\mathbb Z^2$ be arbitrary. Assume
    $0<\epsilon_0t\le B_0\log L$, so these treated frames use supplied
    patches of radius $\epsilon_0t$.
    Let $S$ be the participant set of the block event $R$, put
    $\delta_b=(2b+1)(L^{20})^{-1}$, and assume $b\ge2$; $K$ may be zero.
    Assume bounded point treatments for these data with participants $S$ and
    bounds $(2b,K,2b,\delta_b)$. In the following conclusion each step has
    these same bounds. Write $\mathcal B$ for the repainting baseline,
    $A,B$ for its old and final labels, and $F_M^\lambda(g)$ for its treated
    frame at marks $M$ with homogenizing label $\lambda$.
    Let $E$ be any list of edges and $e$ any edge, with no freshness condition,
    and put $C_e=\mathcal B.\mathrm{nbrLabel}(e)$. Let $X$ be any decidable
    index type, $l$ a list of distinct indices, $x\in l$, and
    $\mathcal A:X\to\mathrm{Frame}$ any family of auxiliary frames.
    The memory is the main frame followed by the auxiliary frames in the order
    $l$. In its initial state the main frame is
    $F_\varnothing^{C_e}(\mathcal B.\mathrm{mainGuide}
        (E,e,\mathrm{edgeStartWord}(A,B,C_e)))$ and the frame at index $x$ is
    $F_\varnothing^{C_e}(\mathcal B.\mathrm{auxGuide}
        (e,\mathrm{auxWord}(A,B,C_e)))$; every other auxiliary frame is
    $\mathcal A(y)$. There is a chain of $9$ source-bounded memory changes to
    the state obtained by replacing those two words by
    $\mathrm{mainWordOne}(A,B,C_e)$ and $\mathrm{auxWordExchanged}(A,B,C_e)$,
    respectively. The main frame stays first, all other auxiliary frames and
    their order are unchanged, and both endpoint sheets are treated at no mark.
    The idle-register layout is empty and its vector is the scalar $1$.
\end{theorem}

\begin{proof}\leanok
    \uses{def:peps_dyadic_edge_band,
        def:peps_dyadic_normal_words,
        def:peps_dyadic_protocol_parties,
        def:peps_dyadic_schedule,
        def:peps_frame,
        def:peps_frame_change,
        def:peps_protocol_bulk_operation_marks,
        def:peps_protocol_local_labels,
        def:peps_protocol_local_point_treatments,
        def:peps_protocol_step_chain,
        def:peps_protocol_treated_frame,
        def:polynomial_peps_model,
        lem:peps_adaptive_patch_supply,
        lem:peps_dyadic_direct_repainting,
        lem:peps_dyadic_unmodified_vertices,
        lem:peps_protocol_change_chains,
        lem:peps_protocol_frame_change_error_mono,
        thm:peps_protocol_bounded_changes,
        thm:peps_protocol_event_segments,
        thm:peps_protocol_lens_exchange_change,
        thm:peps_protocol_local_label_event_uses,
        thm:peps_protocol_small_patch_error_bounds}
    Treat the main sheet ($2$ steps) and $x$-auxiliary sheet ($2$ steps) at the endpoints of $e$, apply the lifted two-sheet lens exchange of Theorem~\ref{thm:peps_protocol_lens_exchange_change} (with error $4L^{-30}\le\delta_b$ by Theorem~\ref{thm:peps_protocol_small_patch_error_bounds} since $b\ge2$), and untreat both sheets ($2+2$ steps) via Lemma~\ref{lem:peps_protocol_change_chains}, giving $2+2+1+2+2=9$ steps.
\end{proof}

\begin{theorem}[The complete edge construction in memory]
    \label{thm:peps_protocol_edge_construction_chain}
    \lean{TNLean.PEPS.Approximation.isStepChain_edgeConstruction}
    \leanok
    \uses{def:peps_adaptive_patch_constraint,
        def:peps_angular_scale,
        def:peps_dyadic_edge_band,
        def:peps_dyadic_protocol_parties,
        def:peps_dyadic_schedule,
        def:peps_frame,
        def:peps_protocol_local_point_treatments,
        def:peps_protocol_repaint_baseline,
        def:peps_protocol_step_chain,
        def:peps_protocol_treated_frame,
        def:polynomial_peps_model,
        lem:peps_adaptive_patch_supply,
        lem:peps_dyadic_unmodified_vertices,
        thm:peps_protocol_constants,
        thm:peps_protocol_event_segments}
    Let $k,m_0,L,b,K,q\in\mathbb N$ with $q\ge1$, and let all the scale
    constants below be real. Fix a family of states and assume the angular-scale
    condition for it with parameters $q,a,A_0,m_0,60,D,L_*$ and $L_*\le L$. Let
    $\Omega$ be a unit vector of the family on the square of side $L$. Assume $m_0\ge2$,
    $A_0\ge1$, the protocol-constant condition on $K_0,\epsilon_0,\nu,a,D,B_0$,
    $L\ge2$ and $2^k<2L$. Set $t=D\log L$. Fix an adaptive patch supply
    $P$ at the cell-center positions for $\Omega$ with parameters
    $B_0,C,c,60,L$, and an arbitrary natural-valued label rank.
    Let $R$ be a repaintable square of hierarchy $k$, with side $n=2^j>K_0t$,
    and let $\mathit{done}\subseteq\mathbb Z^2$ be arbitrary. Assume
    $0<\epsilon_0t\le B_0\log L$, so these treated frames use supplied
    patches of radius $\epsilon_0t$.
    Let $S$ be the participant set of the block event $R$, put
    $\delta_b=(2b+1)(L^{20})^{-1}$, and assume $b\ge2$ and $K\ge1$.
    Assume bounded point treatments for these data with participants $S$ and
    bounds $(2b,K,2b,\delta_b)$. In the following conclusion each step has
    these same bounds. Write $\mathcal B$ for the repainting baseline,
    $A,B$ for its old and final labels, and $F_M^\lambda(g)$ for its treated
    frame at marks $M$ with homogenizing label $\lambda$.
    Let $E$ be a list of edges and $e\notin E$, and put
    $C_e=\mathcal B.\mathrm{nbrLabel}(e)$. Let $X$ be any decidable index
    type, $l$ a list of distinct indices, $x\in l$, and
    $\mathcal A:X\to\mathrm{Frame}$ an arbitrary auxiliary-frame family
    whose frame $\mathcal A(x)$ is the raw frame, all of whose registers are
    held by $\operatorname{labelParty}(R,C_e)$. For arbitrary guide labels
    $\lambda_1,\lambda_2$, the initial memory is
    $F_\varnothing^{\lambda_1}(\mathcal B.\mathrm{completedGuide}(E))$
    followed by $\mathcal A(y)$ in the order $l$.
    There is a chain of $45$ source-bounded memory changes to the memory whose
    main frame is
    $F_\varnothing^{\lambda_2}(\mathcal B.\mathrm{completedGuide}(e::E))$
    and whose auxiliary frame at $x$ is
    $F_\varnothing^{C_e}(\mathcal B.\mathrm{auxGuide}
        (e,\mathrm{auxWordExchanged}(A,B,C_e)))$.
    Every other auxiliary frame and the ordering $l$ are unchanged.
    The idle-register layout is empty with vector $1$; no locality condition
    is imposed on $\lambda_1,\lambda_2$.
\end{theorem}

\begin{proof}\leanok
    \uses{def:peps_dyadic_normal_words,
        def:peps_dyadic_protocol_parties,
        def:peps_dyadic_schedule,
        def:peps_protocol_step_chain,
        def:peps_protocol_treated_frame,
        def:polynomial_peps_model,
        lem:peps_dyadic_unmodified_vertices,
        lem:peps_protocol_change_chains,
        thm:peps_protocol_auxiliary_birth_chain,
        thm:peps_protocol_event_segments,
        thm:peps_protocol_lens_exchange_chain,
        thm:peps_protocol_main_steps_chain,
        thm:peps_protocol_treated_frame_congruence,
        thm:peps_protocol_treated_frame_ends}
    By Theorem~\ref{thm:peps_protocol_treated_frame_ends}, $\mathcal A(x)=F_\varnothing^{C_e}(g_{\mathrm{aux}}(e,\varnothing))$ ($1$ identity step) and $F_\varnothing^{\lambda_1}(g_{\mathrm{done}}(E))=F_\varnothing^{C_e}(g_{\mathrm{main}}(E,e,w_{\mathrm{start}}))$ since $e\notin E$. Concatenating with the lifted $15$-step auxiliary birth chain (Theorem~\ref{thm:peps_protocol_auxiliary_birth_chain}), the $9$-step lens exchange chain (Theorem~\ref{thm:peps_protocol_lens_exchange_chain}), and the lifted $20$-step main-sheet chain (Theorem~\ref{thm:peps_protocol_main_steps_chain}) gives $1+15+9+20=45$ steps.
\end{proof}

\begin{definition}[The final auxiliary frame of an edge]
    \label{def:peps_protocol_auxiliary_final_frame}
    \lean{TNLean.PEPS.Approximation.auxFinal}
    \leanok
    \uses{def:peps_adaptive_patch_constraint,
        def:peps_dyadic_edge_band,
        def:peps_dyadic_protocol_parties,
        def:peps_dyadic_schedule,
        def:peps_frame,
        def:peps_protocol_repaint_baseline,
        def:peps_protocol_treated_frame,
        def:polynomial_peps_model,
        lem:peps_adaptive_patch_supply,
        lem:peps_dyadic_unmodified_vertices}
    Given an adaptive patch supply with parameters $(B_0,C,c,60,L)$, a label rank, a repaintable square $R$ at stage $\mathcal D\subseteq\mathbb Z^2$, and $0<\epsilon_0D\log L\le B_0\log L$, the final auxiliary frame of edge $e$ is $A_e=F_\varnothing^{C_e}(g_{\mathrm{aux}}(e,w_{\mathrm{exch}}))$ (Definition~\ref{def:peps_protocol_treated_frame}), where $C_e$ is the neighbor label across $e$ (\texttt{06-geometry.tex}, lines 198--214 and 620--624).
\end{definition}

\begin{definition}[Auxiliary indices of one event]
    \label{def:peps_protocol_event_auxiliary_indices}
    \lean{TNLean.PEPS.Approximation.eventAuxIndices}
    \leanok
    \uses{def:peps_dyadic_hierarchy,
        def:peps_dyadic_protocol_schedule}
    For a hierarchy depth $k\in\mathbb N$ and a Boolean predicate designating large dyadic squares, the auxiliary-index list of a repainting event $S$ consists of the four pairs $(S,\alpha_e)$ (where $\alpha_e$ denotes the auxiliary initialization of edge $e$) in edge order when $S$ is large, and is empty for small repaintings and junction resizes.
\end{definition}

\begin{theorem}[The auxiliary-index lists]
    \label{thm:peps_protocol_auxiliary_index_facts}
    \lean{TNLean.PEPS.Approximation.filter_largeRepaintingOps,
        TNLean.PEPS.Approximation.auxIndices_eq,
        TNLean.PEPS.Approximation.fst_of_mem_eventAuxIndices,
        TNLean.PEPS.Approximation.nodup_auxIndices,
        TNLean.PEPS.Approximation.mem_auxIndices,
        TNLean.PEPS.Approximation.eq_auxInit_of_mem_auxIndices}
    \leanok
    \uses{def:peps_dyadic_edge_band,
        def:peps_dyadic_hierarchy,
        def:peps_dyadic_protocol_schedule,
        def:peps_protocol_event_auxiliary_indices,
        thm:peps_dyadic_protocol_schedule_counts,
        thm:peps_protocol_from_large_repaintings}
    Filtering a large repainting's operation list for auxiliary initializations yields the four edge initializations in edge order. The global auxiliary-index list is the duplicate-free concatenation of the per-event lists of Definition~\ref{def:peps_protocol_event_auxiliary_indices} along the event schedule, so its elements with first component $S$ are precisely $(S,\alpha_e)$ for the four edges $e$ of a scheduled large repainting $S$.
\end{theorem}

\begin{proof}\leanok
    \uses{def:peps_dyadic_protocol_schedule,
        def:peps_protocol_event_auxiliary_indices,
        thm:peps_dyadic_protocol_schedule_counts,
        thm:peps_protocol_from_large_repaintings}
    By Theorem~\ref{thm:peps_dyadic_protocol_schedule_counts}, each large repainting contributes its $4$ distinct edge initializations and the schedule has no repeated events, so the concatenation is duplicate-free.
\end{proof}

\begin{theorem}[The chain of a large repainting]
    \label{thm:peps_protocol_large_repainting_chain}
    \lean{TNLean.PEPS.Approximation.isStepChain_largeRepainting}
    \leanok
    \uses{def:peps_adaptive_patch_constraint,
        def:peps_angular_scale,
        def:peps_dyadic_edge_band,
        def:peps_dyadic_protocol_parties,
        def:peps_dyadic_schedule,
        def:peps_frame,
        def:peps_protocol_auxiliary_final_frame,
        def:peps_protocol_local_point_treatments,
        def:peps_protocol_repaint_baseline,
        def:peps_protocol_step_chain,
        def:peps_protocol_treated_frame,
        def:polynomial_peps_model,
        lem:peps_adaptive_patch_supply,
        thm:peps_protocol_constants,
        thm:peps_protocol_event_segments}
    Let $k,L,q\in\mathbb N$ with $q\ge1$ and $L\ge2$. Fix real
    $a,A_0,D,L'_0,K_0,\epsilon_0,\nu,B_0,C,c$, $m_0\in\mathbb N$ and a family
    of states. Assume the angular-scale condition for the family with parameters
    $(q,a,A_0,m_0,60,D,L'_0)$, $L'_0\le L$, $m_0\ge4$, $A_0\ge1$,
    and the protocol-constant condition for $(K_0,\epsilon_0,\nu,a,D,B_0)$.
    Let $\Omega$ be a unit vector of the family on the square of side $L$.
    Assume only the upper padding bound $2^k<2L$.

    Fix a patch supply with parameters $B_0,C,c,60,L$ for $\Omega$ at the
    cell centers, an arbitrary rank on guide labels, a repaintable square $R$
    of side $n=2^j>K_0t$, where $t=D\log L$, and an arbitrary stage set
    $\mathit{done}\subseteq\mathbb Z^2$. Supply the radius witness
    $0<\epsilon_0t\le B_0\log L$. Let $b,K\in\mathbb N$ with $b\ge2$ and
    $K\ge1$, and assume bounded point treatments for the parties used by
    the repainting, with bounds
    \begin{align}
        (r,K,r_s,\varepsilon)
         & =\bigl(2b,K,2b,(2b+1)L^{-20}\bigr).
        \label{eq:peps_large_repainting_change_bounds}
    \end{align}
    Let $l$ be a duplicate-free list in $X$, let $e\mapsto x_e\in l$ be injective on the four edges, and let $\mathcal A\colon X\to\mathrm{Frame}$ assign to each $x_e$ the raw frame held by the party of $C_e$. For any homogenizing labels $\lambda_1,\lambda_2$, there is a chain of $189$ source-bounded memory changes with participants in $S_R$ and bounds (\ref{eq:peps_large_repainting_change_bounds}) from the memory with main frame $F_\varnothing^{\lambda_1}(g_{\mathrm{start}})$ and auxiliary frames $\mathcal A(y)$ ($y\in l$) to the memory with completed main frame $F_\varnothing^{\lambda_2}(g_{\mathrm{done}})$, each $\mathcal A(x_e)$ replaced by the final auxiliary frame $A_e$ of Definition~\ref{def:peps_protocol_auxiliary_final_frame}, and all other $\mathcal A(y)$ unchanged (\texttt{06-geometry.tex}, lines 142--221, 328--339 and 572--583).
\end{theorem}

\begin{proof}\leanok
    \uses{def:peps_protocol_auxiliary_final_frame,
        def:peps_protocol_repaint_baseline,
        def:polynomial_peps_model,
        lem:peps_protocol_change_chains,
        thm:peps_protocol_central_birth_chain,
        thm:peps_protocol_edge_construction_chain,
        thm:peps_protocol_treated_frame_congruence}
    Since $\|\Omega\|=1$, the central birth lifts to the memory without an
    additional norm assumption.
    Concatenate its $9$ changes with the $45$ changes on each edge, in the
    order bottom, right, top, left. The distinct auxiliary indices ensure
    that each previous update leaves the next idle copy unchanged. Thus
    $9+4\cdot45=189$, with the same bounds at every step.
\end{proof}

\begin{theorem}[The memory through a large repainting]
    \label{thm:peps_protocol_large_repainting_segment}
    \lean{TNLean.PEPS.Approximation.exists_largeRepainting_segment}
    \leanok
    \uses{def:peps_adaptive_patch_constraint,
        def:peps_angular_scale,
        def:peps_dyadic_edge_band,
        def:peps_dyadic_hierarchy,
        def:peps_dyadic_protocol_parties,
        def:peps_dyadic_protocol_schedule,
        def:peps_frame,
        def:peps_protocol_auxiliary_final_frame,
        def:peps_protocol_local_point_treatments,
        def:peps_protocol_repaint_baseline,
        def:polynomial_peps_model,
        lem:peps_adaptive_patch_supply,
        thm:peps_protocol_constants,
        thm:peps_protocol_event_segments,
        thm:peps_protocol_from_large_repaintings,
        thm:peps_protocol_main_frame_schedule}
    Fix $k,L,q\in\mathbb N$, with $q\ge1$ and $L\ge2$, a family of states,
    and the real angular-scale parameters $(a,A_0,D,L'_0)$ with
    $m_0\in\mathbb N$. Assume the angular-scale condition for the family at
    power $60$, $L'_0\le L$, $m_0\ge4$, $A_0\ge1$, and let $\Omega$ be a
    unit vector of the family on the square of side $L$.
    Fix protocol constants $(K_0,\epsilon_0,\nu,a,D,B_0)$, a patch supply
    with parameters $B_0,C,c,60,L$, a rank on guide labels, and a reference
    repaintable square $R_0$. Assume $L\le2^k<2L$, and explicitly supply
    $t=D\log L>0$ and $\epsilon_0t\le B_0\log L$.

    Let $i$ be less than the event-schedule length, and suppose its $i$th
    event repaints $S$ with side $2^j>K_0t$. Let $b,K\in\mathbb N$ satisfy
    $b\ge2$, $K\ge1$. Whenever $j<k$, assume bounded point
    treatments of $S$ at stage $\mathit{repaintedBlocks}(k,i,j)$, with
    participants used by this event and bounds
    (\ref{eq:peps_large_repainting_change_bounds}). Fix arbitrary frozen
    frames and idle owners for all event-operation pairs. Under the same
    condition $j<k$, require, for every edge of $S$, its frozen frame to equal
    Definition~\ref{def:peps_protocol_auxiliary_final_frame} at that stage,
    and its idle owner to be the party of the label across the edge.
    Then there is a sequence of memory lists indexed by $\mathbb N$ whose
    values at $0$ and $189$ are the protocol states at $i$ and $i+1$.
    For every $s<189$, the transition from $s$ to $s+1$ is a source-bounded
    change with scalar $1$, empty idle-register layout, the same participant set
    and the bounds (\ref{eq:peps_large_repainting_change_bounds})
    (\texttt{06-geometry.tex}, lines 82--100, 142--221, 572--583 and 653--667).
\end{theorem}

\begin{proof}\leanok
    \uses{def:peps_dyadic_protocol_parties,
        def:peps_dyadic_protocol_schedule,
        def:peps_frame,
        thm:peps_protocol_auxiliary_index_facts,
        thm:peps_protocol_constants,
        thm:peps_protocol_from_large_repaintings,
        thm:peps_protocol_large_repainting_chain,
        thm:peps_protocol_main_frame_schedule,
        thm:peps_protocol_treated_frame_ends}
    By Theorems~\ref{thm:peps_protocol_main_frame_schedule}, \ref{thm:peps_protocol_treated_frame_ends}, and~\ref{thm:peps_protocol_auxiliary_index_facts}, the starting and completed unmarked main frames match the main frames at indices $i$ and $i+1$, and the four auxiliary slots $(S,\alpha_e)$ transition from idle raw copies to $A_e$ while all other auxiliary slots remain fixed; Theorem~\ref{thm:peps_protocol_large_repainting_chain} gives the $189$-step chain.
\end{proof}

\begin{definition}[The frozen auxiliary frames]
    \label{def:peps_protocol_frozen_frames}
    \lean{TNLean.PEPS.Approximation.scheduleIndex,
        TNLean.PEPS.Approximation.scheduleIndex_getElem,
        TNLean.PEPS.Approximation.protocolFrozen,
        TNLean.PEPS.Approximation.protocolAuxOwner,
        TNLean.PEPS.Approximation.protocolFrozen_auxInit,
        TNLean.PEPS.Approximation.protocolAuxOwner_auxInit}
    \leanok
    \uses{def:peps_adaptive_patch_constraint,
        def:peps_dyadic_edge_band,
        def:peps_dyadic_hierarchy,
        def:peps_dyadic_protocol_parties,
        def:peps_dyadic_protocol_schedule,
        def:peps_dyadic_schedule,
        def:peps_frame,
        def:peps_protocol_repaint_baseline,
        def:peps_protocol_treated_frame,
        def:polynomial_peps_model,
        lem:peps_adaptive_patch_supply,
        lem:peps_dyadic_unmodified_vertices,
        thm:peps_dyadic_protocol_initial_step,
        thm:peps_protocol_main_frame_schedule}
    In a depth-$k$ hierarchy, the schedule index $i(E)$ of an event $E$ is its first position in the event schedule (or the schedule length if absent), so the $i$th scheduled event $E_i$ satisfies $i(E_i)=i$. Given an adaptive patch supply with parameters $(B_0,C,c,60,L)$, a label rank, and $0<\epsilon_0t\le B_0\log L$, at an auxiliary index $(S,\alpha_e)$ of a level-$j$ repainting $S$ ($j<k$), the frozen frame is $A_e$ (Definition~\ref{def:peps_protocol_auxiliary_final_frame}) and the idle owner is the party of $C_e$ at stage $\mathcal D_{k,i(S),j}$ (and the root raw frame and root party for all other indices) (\texttt{06-geometry.tex}, lines 183--187, 620--624 and 653--667).
\end{definition}

\begin{proof}\leanok
    \uses{thm:peps_dyadic_protocol_schedule_counts}
    Since the event schedule is duplicate-free (Theorem~\ref{thm:peps_dyadic_protocol_schedule_counts}), $i(E_i)=i$, so evaluating at $E_i=S$ recovers $A_e$ and the party of $C_e$ at stage $\mathcal D_{k,i,j}$.
\end{proof}

\begin{theorem}[The holes of the frozen auxiliary frames]
    \label{thm:peps_protocol_frozen_frame_bounds}
    \lean{TNLean.PEPS.Approximation.norm_act_proj_sub_le_of_mem_holes_protocolFrozen,
        TNLean.PEPS.Approximation.length_holes_protocolFrozen_le}
    \leanok
    \uses{def:peps_adaptive_patch_constraint,
        def:peps_dyadic_protocol_parties,
        def:peps_dyadic_protocol_schedule,
        def:peps_frame,
        def:peps_frame_square_patch,
        def:peps_protocol_frozen_frames,
        def:polynomial_peps_model,
        lem:peps_adaptive_patch_supply,
        lem:peps_front_registers,
        thm:peps_protocol_constants,
        thm:peps_protocol_from_large_repaintings}
    In Definition~\ref{def:peps_protocol_frozen_frames} with $L\ge1$ and $0<\epsilon_0t\le B_0\log L$, every hole $h$ of every frozen frame satisfies $\|\Pi_h\Omega-\Omega\|\le L^{-60}$. Under protocol constants with $L\ge2$ and $t=D\log L$, the frozen frame of every auxiliary index for the large-square choice $2^j>K_0t$ has at most $2$ holes (\texttt{06-geometry.tex}, lines 108--118, 620--624 and 653--655).
\end{theorem}

\begin{proof}\leanok
    \uses{def:peps_dyadic_block_guide,
        def:peps_dyadic_edge_band,
        def:peps_dyadic_protocol_parties,
        def:peps_dyadic_schedule,
        def:peps_frame,
        def:peps_protocol_frozen_frames,
        def:peps_protocol_repaint_baseline,
        def:peps_protocol_treated_frame,
        lem:peps_adaptive_patch_supply,
        lem:peps_dyadic_floor,
        lem:peps_dyadic_frozen_hole_count,
        lem:peps_dyadic_unmodified_vertices,
        thm:peps_bulk_steps_treated_frames,
        thm:peps_dyadic_temporary_holes,
        thm:peps_guide_frame_norm,
        thm:peps_protocol_auxiliary_index_facts,
        thm:peps_protocol_bounded_changes,
        thm:peps_protocol_from_large_repaintings,
        thm:peps_protocol_main_frame_schedule,
        thm:peps_treated_frame_spec}
    Each hole uses a supplied patch of radius $\epsilon_0t\le B_0\log L$, giving $\|\Pi_h\Omega-\Omega\|\le L^{-60}$ (Lemma~\ref{lem:peps_adaptive_patch_supply}). By Theorem~\ref{thm:peps_bulk_steps_treated_frames} and Lemma~\ref{lem:peps_dyadic_frozen_hole_count}, the unmarked exchanged auxiliary guide of an edge $e$ has true vertices only at the translates of the two endpoints of $e$, hence at most $2$ holes.
\end{proof}

\begin{theorem}[Enlarging an angular scale]
    \label{thm:peps_angular_scale_monotonicity}
    \lean{TNLean.PEPS.Approximation.IsAngularScale.mono}
    \leanok
    \uses{def:peps_angular_scale}
    For arbitrary $q,m_0,p\in\mathbb N$, real $a,A_0,D,L_0,D'$ and a family of
    states, assume the angular-scale condition for the family with parameters
    $(q,a,A_0,m_0,p,D,L_0)$, $a\ge0$ and $D\le D'$.
    The angular-scale condition then holds with scale $D'$ and threshold
    $\max(L_0,1)$, with all other parameters unchanged. No positivity of
    $A_0$ or $D$ is an additional premise.
\end{theorem}

\begin{proof}\leanok
    \uses{def:peps_angular_scale}
    For $L\ge\max(L_0,1)$, $\log L\ge0$ gives $D\log L\le D'\log L$, and since $a\ge0$ the separation hypotheses at scale $D'$ imply those at scale $D$.
\end{proof}

\begin{theorem}[Dyadic padding of the side length]
    \label{thm:peps_dyadic_padding_bounds}
    \lean{TNLean.PEPS.Approximation.le_two_pow_clog_lt}
    \leanok
    For every integer $L\ge2$, $d=\lceil\log_2L\rceil$ satisfies $d\ge1$ and $L\le2^d<2L$.
\end{theorem}

\begin{proof}\leanok
    \uses{lem:peps_dyadic_padded_side}
    Immediate from $L\ge2$ and Lemma~\ref{lem:peps_dyadic_padded_side}.
\end{proof}

\begin{theorem}[The dyadic distribution protocol from the patch bound]
    \label{thm:peps_protocol_from_angular_and_patch_bounds}
    \lean{TNLean.PEPS.Approximation.largeSizeDyadicProtocolConsequences_of_adaptivePatchConstraint,
        TNLean.PEPS.Approximation.largeSizeDyadicProtocol_of_le_of_adaptivePatchConstraint,
        TNLean.PEPS.Approximation.largeSizeDyadicProtocolBounds_of_adaptivePatchConstraint,
        TNLean.PEPS.Approximation.HierarchySquare.lt_of_eq_pow,
        TNLean.PEPS.Approximation.protocolSquareParty,
        TNLean.PEPS.Approximation.protocolSquareParty_injective,
        TNLean.PEPS.Approximation.exists_eq_inr_of_notMem_range_protocolSquareParty,
        TNLean.PEPS.Approximation.protocolSchedule_fst_eq_of_le,
        TNLean.PEPS.Approximation.near_of_eventUses_repaint}
    \leanok
    \uses{def:peps_adaptive_patch_constraint,
        def:peps_dyadic_protocol}
    Assume the adaptive patch constraint of Proposition~4.1. For every
    $q\in\mathbb N$ with $q\ge2$ and real $J,\Delta>0$, there are a real
    constant $C$ and natural constants $b,\kappa,r,L_0$, chosen before the
    side length and physical instance, such that for every natural
    $L\ge\max(2,L_0)$, square Hamiltonian with parameters $(L,q,J)$, ground
    energy $E_0$ and gapped ground vector $\Omega$ with gap $\Delta$, there
    exist a finite party type and a dyadic distribution protocol with the
    assembly bounds $(C,b,\kappa,r)$ of
    Definition~\ref{def:peps_dyadic_protocol} at $C_N=C_K=C_E=C$; with the patch
    constraint proved, this gives Proposition~\ref{prop:peps_dyadic_protocol}
    (\texttt{06-geometry.tex}, lines 12--53 and 559--681). More precisely, the
    count, expansion and error constants may be taken as in
    Theorem~\ref{thm:peps_dyadic_protocol_explicit_constants}, and $C$ is the
    largest of the three.
\end{theorem}

\begin{proof}\leanok
    \uses{def:peps_dyadic_protocol_parties,
        def:peps_protocol_frozen_frames,
        def:polynomial_peps_model,
        lem:peps_adaptive_patch_supply,
        thm:peps_angular_information_from_gap,
        thm:peps_angular_scale_exists,
        thm:peps_angular_scale_monotonicity,
        thm:peps_dyadic_direct_repainting_constants,
        lem:peps_dyadic_protocol_mesh_ratio,
        thm:peps_dyadic_padding_bounds,
        thm:peps_dyadic_protocol_schedule_counts,
        thm:peps_protocol_constants,
        thm:peps_protocol_from_large_repaintings,
        thm:peps_protocol_frozen_frame_bounds,
        thm:peps_protocol_large_repainting_segment,
        thm:peps_protocol_level_events,
        thm:peps_protocol_point_treatment_change}
    The angular information bound of Proposition~3.4 holds by
    Theorem~\ref{thm:peps_angular_information_from_gap}. Fix the constants in
    the order of the source: $K_0=20$, a small $\epsilon_0$, the mesh ratio
    $\nu$ of Lemma~\ref{lem:peps_dyadic_protocol_mesh_ratio}, below the
    thresholds of the resizes, the point treatments and the direct repaintings, the clearance $a=a_0/2$, the information scale $D$
    of Proposition~3.4 at that clearance, enlarged to at least $1$, and
    $B_0=K_0D$; Proposition~4.1 then supplies the patch constants. Choose
    the count constant $C_N\ge5042$, the expansion constant $C_K\ge G$, the
    prefactor of the small-patch rewrites, the error constant
    $C_E\ge\max\{10080,2b_{\rm pat}+1\}$, where $b_{\rm pat}$ bounds the patches
    of one rewrite, and $\kappa,r,L_0$ from these constants, before $L$ and the
    Hamiltonian. For $L\ge\max(2,L_0)$ take the dyadic depth
    $d=\lceil\log_2L\rceil$ and padded side $N=2^d$, so that $L\le N<2L$.
    Every aligned square of $[0,N]^2$ of side $2^j$, including the root
    square, then has $j\le d$ and block indices below $2^{d-j}$, so it is a
    square of the hierarchy of depth $d$ and its block party is a different
    party for each square, at the scale and address of that square. Every
    other party is the placeholder of a non-root square $S$, placed at $S$;
    each change is recorded as belonging to the repainting of $S$ when its
    scheduled event is that repainting, and the placeholder of $S$ takes part
    only in gates of that repainting. The initial encoded frames are the raw
    copies of $\Omega$, each held by its designated party, and every
    intermediate reference vector is that of the frames of its stage, with
    factors of norm at most $\lVert\Omega\rVert=1$. Every change of the schedule
    keeps the main sheet, so every intermediate list of frames is nonempty. The
    events are listed without repetition
    and each contributes its changes in one block, so the changes of the
    repainting of $S$ form one segment; a party used by that repainting has the
    scale of $S$ or the next one and an address within $6\cdot2^j$ of the
    address of $S$, for $S$ of side $2^j$. The
    large repaintings pass the memory from one state to the next by
    Theorem~\ref{thm:peps_protocol_large_repainting_segment}, with the
    frozen frames and owners of Definition~\ref{def:peps_protocol_frozen_frames}.
    Theorem~\ref{thm:peps_protocol_frozen_frame_bounds} and the schedule
    count bound give at most $2\cdot5040L^3$ holes, and
    Theorem~\ref{thm:peps_protocol_from_large_repaintings} assembles the protocol.
\end{proof}
